\documentclass[UTF8,a4paper,11pt,fontset=fandol]{ctexart}
\usepackage[margin=25mm]{geometry}
\usepackage{amsmath,amssymb,amsthm,mathtools,bm}
\usepackage{booktabs,array,longtable}
\usepackage[colorlinks=true,linkcolor=blue!45!black,citecolor=blue!45!black,urlcolor=blue!45!black]{hyperref}
\usepackage{xcolor}
\usepackage{microtype}
\usepackage{enumitem}
\setlist{nosep}
\allowdisplaybreaks[2]
\setlength{\emergencystretch}{2em}
\newtheorem{historical}{定理}
\renewcommand{\thehistorical}{\Alph{historical}}
\newtheorem{theorem}{定理}
\newtheorem{lemma}{引理}[section]
\newtheorem{proposition}[lemma]{命题}
\newtheorem{corollary}[lemma]{推论}
\theoremstyle{remark}
\newtheorem{remark}[lemma]{注}
\newtheorem{definition}[lemma]{定义}
\renewcommand{\proofname}{证明}
\title{中心立方体极大算子弱 $(1,1)$ 常数的平方根下界}
\author{}
\date{}
\begin{document}
\maketitle
\vspace{-2em}
\begin{abstract}
设 $C_n$ 为 $\mathbb R^n$ 上轴平行立方体中心 Hardy--Littlewood 极大算子的弱 $(1,1)$ 最佳常数。
本文构造非负、有界且紧支撑的测试函数，证明
$\liminf_{n\to\infty}C_n/\sqrt n\ge87/(220\sqrt{58\pi})>1/35$。
构造采用有限域标签排列的周期窄峰及正谱补偿，使一段尺度区间上的归一化方差渐近平坦。
通过节点自身的概率倾斜、有限维高斯比较及带漂移 Brownian 最大值公式，计算局部越阈概率，
再对尺度区间间的交集分别估计。结论沿全部整数维数成立，不要求维数属于特定子序列。
\end{abstract}
\noindent\textbf{关键词：}中心极大算子；弱 $(1,1)$ 不等式；高维立方体；概率倾斜；Brownian 运动
\par\medskip
\section{引言}
对 $f\in L^1(\mathbb R^n)$，定义轴平行立方体的中心 Hardy--Littlewood 极大算子
\[
 M_nf(x)=\sup_{L>0}\frac1{L^n}\int_{x+[-L/2,L/2]^n}|f(y)|\,dy.
\]
本文始终使用 Lebesgue 测度。记 $C_n$ 为其弱 $(1,1)$ 最佳常数，即
\begin{equation}\label{eq:definitionCn}
 C_n=\sup_{\substack{0\ne f\in L^1(\mathbb R^n)\\ \lambda>0}}
 \frac{\lambda\,|\{x:M_nf(x)>\lambda\}|}{\|f\|_1}.
\end{equation}
研究 $C_n$ 随维数的增长，是高维极大函数理论的一个端点问题。
以下三个经典结果给出了相关背景。

\begin{historical}[Stein--Strömberg \cite{SS1983}]
存在绝对常数 $C$，使 $C_n\le Cn\log(e+n)$。
更一般地，此上界对以任意有界对称凸体的同心位似体定义的中心极大算子一致成立。
\end{historical}

\begin{historical}[Aldaz、Aubrun \cite{Aldaz2011,Aubrun2009}]
立方体情形不存在与维数无关的弱 $(1,1)$ 常数。并且对每个 $0<\varepsilon<1$，
存在 $c_\varepsilon>0$，使充分大 $n$ 满足
$C_n\ge c_\varepsilon(\log n)^{1-\varepsilon}$。
\end{historical}

\begin{historical}[Iakovlev--Strömberg \cite{IS2013}]
存在绝对常数 $c>0$，使充分大 $n$ 满足 $C_n\ge c n^{1/4}$。
\end{historical}

另一方面，对每个固定 $1<p<\infty$，立方体极大算子具有与维数无关的
$L^p$ 界\cite{Bourgain2012}。这一强型结论与上述弱端点的维数增长形成对照。
本文关心下界中维数幂次的进一步提高，给出如下结果。

\begin{theorem}\label{thm:main}
轴平行立方体中心极大算子的弱 $(1,1)$ 常数满足
\begin{equation}\label{eq:mainconstant}
 \liminf_{n\to\infty}\frac{C_n}{\sqrt n}
 \ge C_*:=\frac{87}{220\sqrt{58\pi}}.
\end{equation}
这里的极限下确界沿全部正整数维数取值。
\end{theorem}

\begin{theorem}\label{thm:rational}
存在整数 $n_0$，使每个 $n\ge n_0$ 都有
\[
 C_n\ge\frac{\sqrt n}{35}.
\]
\end{theorem}

定理\ref{thm:rational}中的阈值 $n_0$ 未作显式估计。
证明使用有界、非负且紧支撑的 $L^1$ 测试函数，因而不需要将原子测度直接代入
式\eqref{eq:definitionCn}。所得结论是下界，并不确定 $C_n$ 的匹配上界。

构造由一维周期密度的张量积得到。一维密度包含有限域标签确定的窄峰、
正的常数背景和有限 Fourier 补偿项。补偿项使一段尺度区间上的方差在归一化后
渐近平坦；标签计数则同时控制相邻尺度与远离尺度之间的相关性。
在每个短尺度区间内，以首次达到最大值的节点分解越阈事件，并在各节点自身的
概率倾斜下比较有限维高斯分布。随后用带负漂移 Brownian 运动的最大值计算
局部越阈强度。对不同短区间的交集分类求和后，可将这些局部贡献累加。

第\ref{sec:prelim}节建立谱恒等式、标签计数及概率比较所需引理；
第\ref{sec:mainproof}节完成尺度求和、$L^1$ 截断和主定理证明；
第\ref{sec:discussion}节说明常数及结论的适用范围。
可复现数值检验另附报告，不作为本文证明的前提。

\section{预备知识}\label{sec:prelim}
记 $\phi(x)=(2\pi)^{-1/2}e^{-x^2/2}$，$\Phi(x)=\int_{-\infty}^x\phi(t)\,dt$。
$E$ 与 $P$ 分别表示相应概率空间上的期望和概率；更换测度时明确写出下标。
$O(\cdot)$ 和 $o(\cdot)$ 的极限变量默认为 $n\to\infty$；
涉及有限辅助参数时，先固定这些参数。所有有限域下标及周期位移均按相应周期理解。

% 2026-10-13详细展开稿：沿用原标签，供整体替换原01片段；无preamble。
\subsection{正谱补偿与平坦方差}
记 $\mathbb T=\mathbb R/\mathbb Z$，其 Haar 概率测度为 $P$。
以后的一维积分均在 $\mathbb T$ 上进行。取
\[
 s=\frac19,\quad a=\frac{10}{27},\quad\alpha=\frac{729}{29},
 \quad\gamma=\frac{27}{29},\quad I=[10/9,11/9],
\]
\[
 b_m=4\alpha\left[\left(\frac{\sin(\pi ma)}{\pi m}\right)^4
 +3\left(\frac{\sin(\pi ms)}{\pi m}\right)^4\right]\quad(m\ge1).
\]
\begin{lemma}[正谱恒等式]\label{lem:paperplateau}
上述系数非负且 $b_m=O(m^{-4})$。在 $I$ 上，
\[
 F_\infty(L):=1+\sum_{m\ge1}b_m\sin^2(\pi mL)=\gamma L^2,
 \qquad \frac{100}{87}\le\gamma L^2\le\frac{121}{87}.
\]
令 $F_M$ 为截至 $m=M$ 的有限和，则
\[
 0\le F_\infty-F_M\le T_M:=\frac{48}{29M^3},\qquad
 \sqrt{2b_m/3}<\frac{20}{9m^2}.
\]
\end{lemma}
\begin{proof}
\textbf{第一步：计算实线上的三次卷积核。}
记 $T_b(x)=(b-|x|)_+$，$C_b(x)=\int_{\mathbb R}T_b(y)T_b(x-y)dy$。
两因子均为偶函数，故只需 $x\ge0$。当 $0\le x\le b$，交集为
$[x-b,b]$，在 $0,x$ 处分段得到
\begin{align*}
 C_b(x)&=\int_{x-b}^0(b+y)(b-x+y)dy
 +\int_0^x(b-y)(b-x+y)dy\\
 &\qquad+\int_x^b(b-y)(b+x-y)dy.
\end{align*}
第一、三项分别等于
\[
 \int_x^b u(u-x)du=\frac{b^3}{3}-\frac{xb^2}{2}+\frac{x^3}{6},
 \qquad\int_0^{b-x}u(u+x)du
 =\frac{b^3}{3}-\frac{xb^2}{2}+\frac{x^3}{6}.
\]
中间一项等于
\[
 \int_0^x[b(b-x)+xy-y^2]dy=b(b-x)x+\frac{x^3}{6}.
\]
相加即为 $2b^3/3-bx^2+x^3/2$。
当 $b\le x\le2b$，交集为 $[x-b,b]$ 且两绝对值均无需再分段。
令 $c=2b-x$、$u=y-(x-b)$，则
\[
 C_b(x)=\int_{x-b}^b(b-y)(b-x+y)dy
 =\int_0^c(c-u)u\,du=\frac{c^3}{6}.
\]
当 $x\ge2b$，两支撑无内部交集，卷积为0。因此
\[
 C_b(x)=\begin{cases}
 2b^3/3-bx^2+x^3/2,&0\le x\le b,\\
 (2b-x)^3/6,&b\le x\le2b,\\0,&x\ge2b.
 \end{cases}
\]

\textbf{第二步：核对周期化和三次项的抵消。}
将 $T_b$ 作1周期延拓后，它的圆周自卷积等于
$\sum_{j\in\mathbb Z}C_b(x+j)$。所用 $a=10/27<1/2$、$s=1/9$，且
$2a+2s=26/27<1$。当 $x\in[s,2s]$ 时，非零整数平移满足
$|x+j|\ge1-2s=7/9>2a\ge2s$，故没有绕回项；$x=0$ 时同样如此。
在该区间 $x<a$，但 $x\ge s$，所以 $C_a$ 用第一分支，$C_s$ 用第二分支：
\begin{align*}
 C_a(x)+3C_s(x)
 &=\frac{2a^3}{3}-ax^2+\frac{x^3}{2}
   +\frac{(2s-x)^3}{2}\\
 &=\frac{2a^3}{3}+4s^3-6s^2x+(3s-a)x^2.
\end{align*}
因为 $a=3s+3s^2$，最后系数为 $-3s^2$，三次项相消。
令 $C=C_a+3C_s$，于是
\[
 C(0)-C(x)=-2s^3+6s^2x+3s^2x^2.
\]
代入 $s=1/9$、$\alpha=729/29$，常数、一次、二次系数分别为
\[
 1-2\alpha s^3=27/29=\gamma,\qquad
 6\alpha s^2=54/29=2\gamma,\qquad
 3\alpha s^2=27/29=\gamma.
\]
因此 $1+\alpha[C(0)-C(x)]=\gamma(1+x)^2$。

\textbf{第三步：从卷积恒等式到非负谱。}
对非零整数 $m$，直接积分并分部一次得
\begin{align*}
 \widehat T_b(m)&=2\int_0^b(b-x)\cos(2\pi mx)dx\\
 &=\frac{2[1-\cos(2\pi mb)]}{(2\pi m)^2}
 =\left(\frac{\sin(\pi mb)}{\pi m}\right)^2.
\end{align*}
圆周卷积定理给
$\widehat C(m)=\widehat T_a(m)^2+3\widehat T_s(m)^2$。
这两个数均非负。又 $\widehat C(m)=O(m^{-4})$，故 Fourier 级数绝对一致收敛，
可以逐项作差：
\[
 C(0)-C(x)=2\sum_{m\ge1}\widehat C(m)[1-\cos(2\pi mx)]
 =4\sum_{m\ge1}\widehat C(m)\sin^2(\pi mx).
\]
由 $b_m=4\alpha\widehat C(m)$、$x=L-1$ 及
$\sin^2(\pi m(L-1))=\sin^2(\pi mL)$，即得 $F_\infty(L)=\gamma L^2$。
端点值分别为
$(27/29)(10/9)^2=100/87$、$(27/29)(11/9)^2=121/87$。

\textbf{第四步：截断与幅度。}
由 $|\sin|\le1$，
\[
 b_m\le\frac{16\alpha}{\pi^4m^4},\qquad
 \sum_{m>M}m^{-4}\le\int_M^\infty x^{-4}dx=\frac1{3M^3}.
\]
因此尾部至多 $16\alpha/(3\pi^4M^3)<48/(29M^3)$，其中使用 $\pi>3$。
同一界给
\[
 \sqrt{2b_m/3}\le\frac{\sqrt{32\alpha/3}}{\pi^2m^2}
 <\frac{20}{9m^2},
\]
因为 $32\alpha/3=7776/29<400$、$\pi^2>9$。
非负性保证 $F_M\le F_\infty$，遂得全部断言。
\end{proof}

\subsection{确定标签与非零状态计数}
\begin{lemma}[有限域标签的存在]\label{lem:papersidon}
令 $B=2^b$，$b\ge2$，$q=2^d$，在 $\mathbb F_q$ 中固定本原元 $g$。
置 $N_B=B-1$、$T_B=N_B(N_B+1)/2$、$R_B=N_B(N_B-1)$ 和
\[
 H(B)=N_B+T_B(T_B-1)+2R_B^2+4R_B<3B^4.
\]
若 $q>H(B)$，存在 $\mathbb F_2$ 独立的 $\beta_0,\ldots,\beta_{b-1}$，
使其线性张成空间 $W$ 的非零元素构成乘法 Sidon 集。即两个非零元素
的乘积相等只可能来自相同无序多重集。令 $K=q-1$，则
\[
 \mathcal D=\{\log_g(a/c):a,c\in W\setminus\{0\},\ a\ne c\}
\]
有 $R_B$ 个元素，$\mathcal D\cup\{0\}$ 在 $\mathbb Z/K\mathbb Z$ 中无邻接，
且 $\pm1,\pm2\notin\mathcal D$。
\end{lemma}
\begin{proof}
\textbf{第一步：有限域多项式的根数预算。}
若 $P$ 是总次数为 $r$ 的非零多项式，独立均匀取其各变量于 $\mathbb F_q$，
则 $\Pr(P=0)\le r/q$。单变量情形来自根数至多次数。
对变量数归纳，将 $P$ 看成末变量 $X_b$ 的 $j$ 次多项式，首项系数
$P_j$ 非零、总次数至多 $r-j$。在 $P_j=0$ 的纤维，全部概率不超过
$(r-j)/q$；在其余纤维，末变量至多有 $j$ 个根，再付 $j/q$。
两者相加给所述预算；特征2不影响这一根数论证。

\textbf{第二步：线性独立及乘法 Sidon 条件。}
独立取 $X_0,\ldots,X_{b-1}\in\mathbb F_q$，对每个非零
$a\in\mathbb F_2^b$ 定义 $L_a(X)=\sum_i a_iX_i$。
共有 $N_B=B-1$ 个非零一次式，故某个 $L_a=0$ 的概率至多 $N_B/q$。
这些零事件均不发生，恰保证所选 $X_i$ 对 $\mathbb F_2$ 线性独立。

在 $N_B$ 个非零二进制向量上，无序二元素多重集有
$T_B=N_B(N_B+1)/2$ 个。不同多重集 $\{a,c\}\ne\{e,f\}$ 的
多项式 $L_aL_c-L_eL_f$ 不恒为0。确切地，多项式环
$\mathbb F_q[X_0,\ldots,X_{b-1}]$ 是唯一分解整环，每个非零线性式
不可约。若两乘积恒等，则一侧的线性因子必与另一侧某因子成比例。
两个二进制非零线性式成比例时，选某个系数为1的位置可见比例必为1，
故相应二进制向量相等；消去因子后另一向量也相等。
这与多重集不同矛盾。重复因子 $L_a^2$ 在特征2下仍是非零多项式，
故同样包含在论证中。
不同多重集对有 $\binom{T_B}{2}$ 个，每个二次式付 $2/q$，总误差为
$T_B(T_B-1)/q$。

在上述两类坏事件的补集，令 $\beta_i=X_i$，则 $W$ 有 $B$ 个元素，
$W\setminus\{0\}$ 满足乘法 Sidon 条件。
若两个非平凡有序比率 $a/c=e/f$，则 $af=ce$。
Sidon 性使 $\{a,f\}=\{c,e\}$。其中 $a=c,f=e$ 与 $a\ne c$ 矛盾，
故只能 $a=e,c=f$。所以全部 $R_B=N_B(N_B-1)$ 个非平凡有序比率互异。

\textbf{第三步：排除邻接比率及短 位移。}
对任意两个非平凡有序比率 $L_a/L_c$、$L_f/L_e$，邻接关系对应
$L_aL_e=gL_cL_f$。多项式 $L_aL_e-gL_cL_f$ 不恒为0：
第二乘积有某个非零单项式系数，等于1；差式中该位置的系数是
$-g$ 或 $1-g$，均非零，因为本原元 $g\ne0,1$。
共 $R_B^2$ 个二次式，误差至多 $2R_B^2/q$。
故在补集上 $\mathcal D\cap(\mathcal D+1)=\varnothing$。

此外，对每个非平凡有序对 $a,c$ 和 $j\in\{1,-1,2,-2\}$，排除
$L_a=g^jL_c$。由 $g^j\ne1$，此一次式非零，共 $4R_B$ 个式子，
误差至多 $4R_B/q$。于是 $\pm1,\pm2\notin\mathcal D$。
比率1本来已由 $a\ne c$ 排除；再结合 $\pm1\notin\mathcal D$，
得到 $\mathcal D\cup\{0\}$ 无邻接。

\textbf{第四步：合并存在性阈值。}
联合坏概率至多 $H(B)/q<1$，因此至少有一组同时合格的选择。
又 $T_B<B^2/2$、$R_B<B^2$、$N_B<B$，且 $B\ge4$，所以
\[
 H(B)<\frac94B^4+4B^2+B
 \le\left(\frac94+\frac14+\frac1{64}\right)B^4
 =\frac{161}{64}B^4<3B^4.
\]
每个 $q=2^d$ 都有有限域：在 $X^q-X$ 的分裂域中，其根集对加法、乘法
及非零逆元封闭，导数为 $-1$ 保证有 $q$ 个互异根。
最后说明本原元的存在。有限乘法群是交换群；把所有元素的阶按各素数
取最高幂，选实现这些幂的元素并相乘，可得到阶等于群指数 $e$ 的元素。
群中每个元素都是 $X^e-1$ 的根，故 $q-1\le e$；另一方面元素阶整除
群阶，故 $e\le q-1$。于是 $e=q-1$，所得元素为本原元。
\end{proof}

取 $x_j=g^j$，$0\le j<K$，令
\[
 y_i(j)=\operatorname{Tr}(\beta_ix_j),\quad
 y_j=\sum_{i=0}^{b-1}2^iy_i(j),\quad
 U_j=(y_j+1/2)/B,\quad a_j=(j+U_j)/K.
\]
为明确标签的联合计数，先写出对应线性矩阵。
迹映射是
\[
 \operatorname{Tr}(x)=x+x^2+\cdots+x^{2^{d-1}}\in\mathbb F_2.
\]
该多项式非零且次数 $q/2<q$，因而不可能在全域恒为0。
迹只取0或1，故存在迹为1的元素。对任意 $c\ne0$，乘法 $x\mapsto cx$
遍历全域，所以 $x\mapsto\operatorname{Tr}(cx)$ 也是非零线性泛函。
因此迹配对将域元素 $c$ 与线性泛函 $\operatorname{Tr}(c\,\cdot)$ 一一对应。

对 位移 $k$，把 $\mathbb F_q$ 看作 $d$ 维 $\mathbb F_2$ 向量空间。
联合标签映射的 $2b$ 行分别对应
\[
 \beta_0,\ldots,\beta_{b-1},g^k\beta_0,\ldots,g^k\beta_{b-1}.
\]
非退化迹配对保证该矩阵的秩为
$\dim(W+g^kW)=2b-\dim(W\cap g^kW)$。
若 $k\notin\mathcal D\cup\{0\}$，交集只有0，矩阵秩为 $2b$；
每个输出 $(y,z)$ 的纤维有 $2^{d-2b}=q/B^2$ 个状态。
若 $k=0$，秩为 $b$，且输出必须满足 $y=z$，每个对角纤维有 $q/B$ 个状态。
若 $k\in\mathcal D$，Sidon 性保证交集仅有一个非零元素：
两个不同非零交点会给同一个比率 $g^k$ 的两种表示，与有序比率互异矛盾。
故交集维数为1，联合矩阵秩为 $2b-1$。

用二进制向量 $a,c$ 表示这个匹配，即
$\beta(a):=\sum_i a_i\beta_i=g^k\sum_i c_i\beta_i=:g^k\beta(c)$。
匹配是唯一的，且 $a,c\ne0$、$a\ne c$；若 $a=c$，便有
$(1-g^k)\beta(a)=0$，迫使 $k=0$。
令标签字符 $\chi_a(y)=(-1)^{\sum_i a_iy_i}$。
全域均匀状态下，缺陷 位移 的联合输出律恰为
$B^{-2}[1+\chi_a(y)\chi_c(z)]$。
这是因为其像为唯一线性关系 $a\cdot y+c\cdot z=0$ 的超平面，
每个像内点纤维大小为 $q/2^{2b-1}$。

实际序列 $x_j=g^j$ 遍历的是全部非零状态，不含0。
因此只须从全域输出的全零标签纤维减去一次，而不能改动其余纤维：
\begin{equation}\label{eq:paperlabelcount}
 N_k(y,z)=q/B^2-\mathbf1_{\{y=z=0\}}
 \quad(k\notin\mathcal D\cup\{0\}).
\end{equation}
零位移 的计数为
$qB^{-1}\mathbf1_{\{y=z\}}-\mathbf1_{\{y=z=0\}}$；
缺陷 位移 的计数为
$qB^{-2}[1+\chi_a(y)\chi_c(z)]-\mathbf1_{\{y=z=0\}}$。
这些公式均为确定性矩阵秩与纤维计数，不使用随机独立标签假设。

定义分格原函数
\[
 \mathcal P((j+r)/K)=\mathbf1_{\{r\ge U_j\}}-r\quad(0\le r<1).
\]
它在格端点为0，周期意义下 $\mathcal P'=\sum_j\delta_{a_j}-K$。
置 $H_L(x)=\mathcal P(x+L/2)-\mathcal P(x-L/2)$，则 $\int H_L=0$、$|H_L|\le2$。
令 $\beta_\delta=\delta^{-1}\mathbf1_{[-\delta/2,\delta/2]}$ 周期延拓，
$H_{L,\delta}=H_L*\beta_\delta$，并置
\[
 B_L(x)=\sum_{m\le M}\sqrt{2b_m/3}\sin(\pi mL)\cos(2\pi mx),
 \quad G_\delta(L,J)=\int(H_{L,\delta}+B_L)(H_{J,\delta}+B_J).
\]

\begin{lemma}[交叉项与平滑误差]\label{lem:papercross}
令 $a_B=1-1/B$，$1\le M<K$。对上述确定标签源，全部 $L,J$ 有
\[
 \left|\int(H_{L,\delta}B_J+B_LH_{J,\delta})\right|\le C_x,
\]
\[
 C_x=\frac{40}{9}\left[
 \frac{a_B\sqrt{K+1}}K\sum_{m\le M}m^{-2}
 +\frac{\pi a_B^2}{2K}\sum_{m\le M}m^{-1}\right].
\]
若 $\delta\le1/K$，则全部原子项的二矩的平滑误差至多 $2K\delta/3$。
\end{lemma}
\begin{proof}
\textbf{第一步：Gauss 和的模。}
记 $\eta_m(g^j)=e^{2\pi imj/K}$、$\chi_c(x)=(-1)^{\operatorname{Tr}(cx)}$。
因为 $1\le m<K$，$\eta_m$ 是非平凡乘法字符。
乘以某个使 $\eta_m(t_0)\ne1$ 的元素重排求和，便得
$\sum_{t\ne0}\eta_m(t)=0$。
对 $c\ne0$，迹配对非退化，故 $\chi_c$ 是非平凡加法字符；
平移一个使 $\chi_c(x_0)=-1$ 的元素重排全域求和，得到
$\sum_x\chi_c(x)=0$，从而 $\sum_{x\ne0}\chi_c(x)=-1$。

设 $G_{c,m}=\sum_{x\ne0}\chi_c(x)\eta_m(x)$。
将其模平方展开，利用字符的乘法性，作 $x=ty$ 变量替换得
\begin{align*}
 |G_{c,m}|^2
 &=\sum_{x,y\ne0}\chi_c(x-y)\eta_m(x/y)\\
 &=\sum_{t\ne0}\eta_m(t)\sum_{y\ne0}\chi_c((t-1)y)\\
 &=K-\sum_{t\ne0,1}\eta_m(t)=K+1=q.
\end{align*}
$t=1$ 的内和为 $K$，其余内和均为 $-1$。
此论证只需要 $m\not\equiv0\pmod K$，不需要 $m$ 与 $K$ 互素。

\textbf{第二步：保留格内相位的实际原子系数。}
由 $y_i=(1-\chi_{\beta_i}(x_j))/2$ 和 $\sum_i2^i=B-1$，
\begin{align*}
 V_j:=U_j-1/2
 &=\frac1B\left[\sum_i2^i\frac{1-\chi_{\beta_i}(x_j)}2
 +\frac12-\frac B2\right]\\
 &=-\sum_i\frac{2^i}{2B}\chi_{\beta_i}(x_j).
\end{align*}
绝对系数和为 $(B-1)/(2B)=a_B/2$，故 $|V_j|\le a_B/2$；
第一步的 Gauss 和又给
\[
 \left|\sum_jV_je^{2\pi imj/K}\right|
 \le\frac{a_B\sqrt q}{2}.
\]
令 $t=2\pi m/K$，则实际原子 Fourier 和为
\[
 Z_m=\sum_je^{2\pi im(j+U_j)/K}
 =e^{it/2}\sum_je^{2\pi imj/K}e^{itV_j}.
\]
准确保留 $V_j$ 相位，用积分余项
$e^{iz}=1+iz+R(z)$、$|R(z)|\le z^2/2$。
常数项因 $\sum_je^{2\pi imj/K}=0$ 消失，于是
\[
 |Z_m|\le t\frac{a_B\sqrt q}{2}+K\frac{t^2a_B^2}{8}.
\]
除以 $\pi m$ 并代入 $t$，得
\[
 |\zeta_m|:=\frac{|\Re Z_m|}{\pi m}
 \le\frac{a_B\sqrt q}{K}+\frac{\pi m a_B^2}{2K}.
\]
这里没有删除格内相位，也没有使用零频或 $m=K$ 的混叠频率。

\textbf{第三步：原函数 Fourier 系数与两向交叉。}
在分布意义下 $\mathcal P'=\sum_j\delta_{a_j}-K$。
用周期分部积分，非零频率 $m$ 满足
\[
 \Re Z_m=\int\cos(2\pi mx)\,d\mathcal P'(x)
 =2\pi m\int\mathcal P(x)\sin(2\pi mx)dx.
\]
因而 $\int\mathcal P\sin(2\pi mx)=\Re Z_m/(2\pi m)$。
对平移差 $H_L$，改变积分变量并展开余弦，得到
\begin{align*}
 \int H_L(x)\cos(2\pi mx)dx
 &=\int\mathcal P(y)[\cos(2\pi my-\pi mL)
 -\cos(2\pi my+\pi mL)]dy\\
 &=2\sin(\pi mL)\int\mathcal P(y)\sin(2\pi my)dy
 =\zeta_m\sin(\pi mL).
\end{align*}
令 $c_m=\sqrt{2b_m/3}$，则两向交叉恰为
\[
 \int(H_LB_J+B_LH_J)
 =2\sum_{m\le M}c_m\zeta_m\sin(\pi mL)\sin(\pi mJ).
\]
箱平滑算子在余弦频率 $m$ 上的乘子为
\[
 \int_{-\delta/2}^{\delta/2}\frac{\cos(2\pi my)}\delta dy
 =\frac{\sin(\pi m\delta)}{\pi m\delta},
\]
其模不超过1。因此平滑后的同一交叉式只增加此乘子。
逐项使用 $2c_m<40/(9m^2)$、$|\sin|\le1$ 及第二步的 $|\zeta_m|$ 界，
即得 $C_x$，整个估计对所有 $L,J$ 一致。

\textbf{第四步：移动窗口的平方增量。}
先取 $0\le z\le\delta\le1/K$。由原函数的导数，
\[
 H_L(x+z)-H_L(x)=N_+(x)-N_-(x),
\]
其中 $N_+$、$N_-$ 分别计数两个端点的长度 $z$ 移动窗口中的原子。
原函数每个窗口内的常数斜率贡献均为 $-Kz$，在相减中相消。
每格恰有一个原子；长度不超过 $1/K$ 的窗口至多与两个格的内部相交，
所以 $N_\pm\in\{0,1,2\}$。这不要求原子之间有额外的最小间距。
于是
\[
 (N_+-N_-)^2\le N_+^2+N_-^2\le2(N_++N_-).
\]
对每个固定原子，包含它的窗口中心集合长度为 $z$，故
$\int N_+=\int N_-=Kz$，从而
$\|H_L(\cdot+z)-H_L\|_2^2\le4Kz$。
负 $z$ 由平移不变性得到同一 $4K|z|$ 界。

\textbf{第五步：从对角能量到任意矩阵元。}
令 $S$ 为宽度 $\delta$ 的箱平滑。它自伴且是 $L^2$ 收缩，故
$A=I-S^2$ 半正定。$S^2$ 的平移核为
$\nu_\delta(z)=(\delta-|z|)_+/\delta^2$，且
\[
 \int|z|\nu_\delta(z)dz
 =\frac2{\delta^2}\int_0^\delta z(\delta-z)dz=\frac\delta3.
\]
用平移不变性和偶核展开平方得
\begin{align*}
 \langle H_L,AH_L\rangle
 &=\frac12\int\|H_L(\cdot+z)-H_L\|_2^2\nu_\delta(z)dz\\
 &\le\frac12\int4K|z|\nu_\delta(z)dz=\frac23K\delta.
\end{align*}
半正定型的 Cauchy 不等式给
\[
 |\langle H_L,AH_J\rangle|
 \le\langle H_L,AH_L\rangle^{1/2}
 \langle H_J,AH_J\rangle^{1/2}\le\frac23K\delta.
\]
又 $\langle H_L,AH_J\rangle=\langle H_L,H_J\rangle-
\langle SH_L,SH_J\rangle$，故得到全部原子项的二矩的平滑误差。
\end{proof}

\subsection{协方差核估计}
\begin{lemma}[连续安全尺度上的核]\label{lem:paperkernel}
令
\[
 \mathcal S=\{L\in I:\operatorname{dist}(K(L-1)\bmod K,\mathcal D)>2\},
\]
\[
 \quad\eta_D=\frac1{3B^2}+\frac2K,\quad\eta_N=\frac2{3B^2}+\frac2K,
 \quad e=C_x+\frac{2K\delta}{3}.
\]
在 $\mathcal S$ 上，
\[
 V_{\min}:=100/261-T_M/3-\eta_D-e\le G_\delta(L,L),
 \quad\left|\frac{G_\delta(L,L)}{\gamma L^2/3}-1\right|
 \le r_*:=\frac{T_M+3(\eta_D+e)}{100/87}.
\]
若 $V_{\min}\ge1/4$，记 $\tau=K|L-J|/2$ 和
$\kappa=G_\delta(L,J)/\sqrt{G_\delta(L,L)G_\delta(J,J)}$，则
\[
 1-\kappa\ge\frac{174}{121}\tau-8(\eta_N+e)\quad(0<\tau\le1/8),
\]
\[
 1-\kappa\ge\frac{14703}{61952}-8(\eta_N+e)\quad(1/8\le\tau\le1).
\]
沿 $B,K,M\to\infty$、$e\to0$，全部安全尺度最终有
$\kappa\ge-1/2$，而 $1/8\le\tau\le K/18$ 时 $1-\kappa\ge1/5$。
\end{lemma}
\begin{proof}
\textbf{第一步：量化中点的误差及其原函数。}
对标签字符 $\chi_a$，定义
\[
 q_y(r)=\mathbf1_{\{r\ge(y+1/2)/B\}}-r,
 \quad f_a(r)=B^{-1}\sum_{y=0}^{B-1}\chi_a(y)q_y(r).
\]
零字符给 $f_0=e_B$，其中
$e_B(r)=B^{-1}\#\{y:(y+1/2)/B\le r\}-r$。
在第 $j$ 个量化格内写 $r=(j+t)/B$，$0\le t\le1$，则
\[
 e_B((j+t)/B)=\begin{cases}-t/B,&0\le t<1/2,\\(1-t)/B,&1/2\le t\le1.
 \end{cases}
\]
特别地 $|e_B|\le1/(2B)$，每格积分为0，而且
\[
 \int_0^1e_B(r)^2dr
 =B\cdot B^{-3}\left[\int_0^{1/2}t^2dt+
 \int_{1/2}^1(1-t)^2dt\right]=\frac1{12B^2}.
\]
置 $E_B(x)=\int_0^xe_B(r)dr$。每格积分为0意味着 $E_B(j/B)=0$，
而格内准确公式为
\[
 E_B((j+t)/B)=\begin{cases}-t^2/(2B^2),&0\le t\le1/2,\\
 -(1-t)^2/(2B^2),&1/2\le t\le1.
 \end{cases}
\]
因此 $-1/(8B^2)\le E_B\le0$，且
\[
 \int_0^1E_B(r)dr=-\frac1{24B^2}.
\]
这些界针对整个连续区间，不是离散节点抽样。

\textbf{第二步：全部非零 Walsh 字符的能量。}
当 $a\ne0$，$\sum_y\chi_a(y)=0$，故 $f_a$ 中的 $-r$ 项消失；
$f_a(r)$ 正是带符号前缀和除以 $B$。
令 $\ell$ 为 $a$ 的最低置位，$q_a=2^\ell$。
把标签顺序分成 $2q_a$ 个标签一块。在每块中，最低置位之前的字符均为0，
所以前 $q_a$ 个字符具有同一个符号，后 $q_a$ 个为相反符号；
更高位只决定该块整体的正负号。因此块内前缀的绝对值依次上升到
$q_a$，再下降到0，块与块之间没有累积残差。
每个中点阈值相隔 $1/B$，块端两段各为半步但前缀值均为0。
各值 $j=1,\ldots,q_a-1$ 出现两段、每段长 $1/B$，最高值 $q_a$
出现一段长 $1/B$。每块平方积分遂为
\[
 \frac1{B^3}\left[2\sum_{j=1}^{q_a-1}j^2+q_a^2\right]
 =\frac1{B^3}\left[\frac{(q_a-1)q_a(2q_a-1)}3+q_a^2\right]
 =\frac{2q_a^3+q_a}{3B^3}.
\]
全区间共有 $B/(2q_a)$ 块，故
\begin{equation}\label{eq:paperwalshenergy}
 \|f_a\|_2^2=\frac{q_a^2}{3B^2}+\frac1{6B^2}.
\end{equation}
$q_a=B/2$ 只可能由唯一字符 $a=2^{b-1}$ 产生，其能量为
$1/12+1/(6B^2)$。其他字符均 $q_a\le B/4$，能量至多
$1/48+1/(6B^2)$。
因此任意两个不同非零字符 $a,c$ 的任何截段平移内积均满足
\[
 |\int_J f_a(r)f_c(r+\theta)dr|
 \le\|f_a\|_2\|f_c\|_2
 \le Q_B:=\sqrt{(1/12+1/(6B^2))(1/48+1/(6B^2))}.
\]
变量平移若带一次周期换格，同样由相应子区间的 $L^2$ 范数控制。
至少一个字符不是唯一最高层，故此处得到 $Q_B=1/24+O(B^{-2})$，
而不是使用两者都取最高能量的较粗界。

\textbf{第三步：实际自相关的 Walsh 逆变换与零状态校正。}
记 $C(z)=\int\mathcal P(x)\mathcal P(x+z/K)dx$，写 $z=k+\theta$，
$0\le\theta\le1$。当 $x=(j+r)/K$、$r<1-\theta$ 时，平移后在第
$j+k$ 格；当 $r\ge1-\theta$ 时，在第 $j+k+1$ 格。
定义
\[
 I^1_{yw}(\theta)=\int_0^{1-\theta}q_y(r)q_w(r+\theta)dr,
 \quad I^2_{yw}(\theta)=\int_{1-\theta}^1q_y(r)q_w(r+\theta-1)dr.
\]
换元 $dx=dr/K$ 后，准确计数式为
\[
 C(z)=K^{-1}\sum_{y,w}[N_k(y,w)I^1_{yw}+N_{k+1}(y,w)I^2_{yw}].
\]
令 $W^\nu_{ac}=B^{-2}\sum_{y,w}\chi_a(y)\chi_c(w)I^\nu_{yw}$。
对全域均匀状态，字符平均满足
\[
 E[\chi_a(y(x))\chi_c(y(g^kx))]
 =\begin{cases}1,&\beta(a)=g^k\beta(c),\\0,&\text{否则}.
 \end{cases}
\]
这是对加法字符 $(-1)^{\operatorname{Tr}((\beta(a)+g^k\beta(c))x)}$
的全域平均。因而 Walsh 逆变换只保留零字符和匹配对。
令 $\mathcal M_k=\{(a,c):a,c\ne0,\ \beta(a)=g^k\beta(c)\}$，则
\begin{align*}
 C(z)&=\frac qK C_\infty(z)-\frac{i(\theta)}K,\\
 C_\infty(z)&=W^1_{00}+W^2_{00}
 +\sum_{\mathcal M_k}W^1_{ac}+\sum_{\mathcal M_{k+1}}W^2_{ac},\\
 i(\theta)&=I^1_{00}+I^2_{00}.
\end{align*}
这里下标 $00$ 在 $I$ 中表示实际全零标签，在 $W$ 中表示零字符，含义不同。
其来源分别是从实际计数减去状态0，以及全域均匀标签律的常数项。
因为 $q=K+1$，这也可写成
\[
 C(z)=C_\infty(z)+[C_\infty(z)-i(\theta)]/K.
\]

为控制最后的修正，注意每个全域联合标签律均非负，且两边际均匀。
将两段合并后，对联合状态和空间变量用 Cauchy，得到
\[
 |C_\infty(z)|\le B^{-1}\sum_y\int_0^1q_y(r)^2dr.
\]
确切地 $B^{-1}\sum_yq_y(r)^2=r(1-r)+(1-2r)e_B(r)$。
分部积分时 $E_B(0)=E_B(1)=0$，因此
\[
 B^{-1}\sum_y\int q_y^2
 =\frac16+\int_0^1(1-2r)E_B'(r)dr
 =\frac16+2\int_0^1E_B(r)dr
 =\frac16-\frac1{12B^2}\le\frac16.
\]
此外 $i(\theta)$ 是 $q_0$ 的周期自相关，模至多
\[
 \|q_0\|_2^2=\int_0^{u_0}r^2dr+\int_{u_0}^1(1-r)^2dr
 =u_0^2-u_0+1/3\le1/3,\qquad u_0=1/(2B).
\]
故 $|C(z)-C_\infty(z)|\le1/(2K)$。

若 $k,k+1$ 均非零位移，无邻接条件使两段中至多一个非零匹配，
且它的两个字符不同。另一方面
\[
 W^1_{00}+W^2_{00}=\int_0^1e_B(r)e_B(r+\theta\bmod1)dr,
\]
其模至多 $1/(12B^2)$。第二步的异字符界遂给
\begin{equation}\label{eq:paperfarC}
 |C(z)|\le Q_B+\frac1{12B^2}+\frac1{2K}=:\mathcal C_{BK}.
\end{equation}
若两段均非缺陷，则没有异字符项，故可使用更小的
$|C(z)|\le1/(12B^2)+1/(2K)$。

\textbf{第四步：靠近零位移的精确量化误差。}
取 $0\le\tau\le1$。其第一段 位移 为0，第二段为1；$1\notin\mathcal D$
保证第二段标签联合律为满秩均匀律。第一段的同标签平均是
\begin{align*}
 B^{-1}\sum_yq_y(r)q_y(r+\tau)
 &=F_B(r)-rF_B(r+\tau)-(r+\tau)F_B(r)+r(r+\tau)\\
 &=r(1-r-\tau)+(1-r-\tau)e_B(r)-r e_B(r+\tau),
\end{align*}
其中 $F_B(r)=r+e_B(r)$，因为 $r\le r+\tau$ 时两指示函数的乘积
等于较左指示函数。令 $a_0=1-\tau$，两误差项逐项分部积分：
\[
 \int_0^{a_0}(a_0-r)e_B(r)dr=\int_0^{a_0}E_B(r)dr,
\]
\[
 -\int_0^{a_0}r e_B(r+\tau)dr
 =-a_0E_B(1)+\int_0^{a_0}E_B(r+\tau)dr
 =\int_0^{a_0}E_B(r+\tau)dr.
\]
故第一段总量化误差恰为
$\int_0^{a_0}[E_B(r)+E_B(r+\tau)]dr$，模至多 $a_0/(4B^2)$。
第二段是 $e_B(r)e_B(r+\tau-1)$ 的积分，区间长 $\tau$，模至多
$\tau/(4B^2)$。连续均匀标签的第一段积分为
\[
 \int_0^{a_0}r(a_0-r)dr=\frac{a_0^3}{6}.
\]
加入第三步的实际非零状态校正，即得
\[
 |C(\tau)-(1-\tau)^3/6|\le\frac1{4B^2}+\frac1{2K}.
\]

\textbf{第五步：反射项、对角和完整原子项的协方差核。}
自相关 $C$ 是偶函数、$K$ 周期函数。将 $H_LH_J$ 的四项展开，
两个同向平移给相同的差位移，两反向平移给相同的和位移，因此
\[
 HH_{LJ}:=\int H_LH_J
 =2C(K(L-J)/2)-2C(K(L+J)/2).
\]
不妨 $L\le J$，令反射参数
$z_r=K((L+J)/2-1)\in[K/9,2K/9]$。它到 $K(L-1)$ 的距离为
$K(J-L)/2=\tau$。安全条件给 $\operatorname{dist}(K(L-1),\mathcal D)>2$，
所以在 $\tau\le1$ 时 $z_r$ 到缺陷的距离大于1，floor 和 next 单元
都不是缺陷；又 $z_r$ 在上述区间内，两个 单元 都非零位移。
第三步的满秩小界于是控制反射项，得到
\[
 |HH_{LJ}-(1-\tau)^3/3|
 \le2\left(\frac1{4B^2}+\frac1{2K}\right)
 +2\left(\frac1{12B^2}+\frac1{2K}\right)=\eta_N.
\]
对角时，$C_\infty(0)=1/6-1/(12B^2)$ 是第三步已算出的精确值；
反射仍用同一满秩界，故
\[
 |HH_{LL}-1/3|
 \le\frac1{6B^2}+\frac1{6B^2}+\frac2K=\eta_D.
\]
这些界已经包含反射和实际零状态校正。

\textbf{第六步：低谱、完整矩及方差比例。}
余弦正交性 $\int\cos(2\pi mx)\cos(2\pi\ell x)dx=
\mathbf1_{m=\ell}/2$ 给
\[
 BB_{LJ}=\frac13\sum_{m\le M}b_m\sin(\pi mL)\sin(\pi mJ),
 \qquad BB_{LL}=\frac{F_M(L)-1}{3}.
\]
由 $F_M(L)\le121/87$ 和 Cauchy，$|BB_{LJ}|\le34/261$，故上界为此数。
下界可以更小：写 $L=1+x$，$x\in[1/9,2/9]$，则
$\sin(\pi mL)=(-1)^m\sin(\pi mx)$。
对 $m=1,2,3,4$，对应角区间分别为
$[\pi/9,2\pi/9]$、$[2\pi/9,4\pi/9]$、$[\pi/3,2\pi/3]$、
$[4\pi/9,8\pi/9]$，正弦均为正。因此每个前四频的两尺度乘积非负，
其余频率按尾部付费，得到
\[
 BB_{LJ}\ge-\frac13\sum_{m>4}b_m\ge-\frac{T_4}{3}=-\frac1{116}.
\]
若 $M<4$，所有保留项本来非负，同一下界仍成立。

完整 Gram 的交叉和平滑误差至多 $e=C_x+2K\delta/3$。
于是 $|G_\delta(L,L)-F_M(L)/3|\le\eta_D+e$。
由 $F_M(L)\ge\gamma L^2-T_M$ 得
\[
 G_\delta(L,L)\ge100/261-T_M/3-\eta_D-e=V_{\min},
\]
\[
 \left|\frac{G_\delta(L,L)}{\gamma L^2/3}-1\right|
 \le\frac{T_M+3(\eta_D+e)}{\gamma L^2}
 \le\frac{T_M+3(\eta_D+e)}{100/87}=r_*.
\]

\textbf{第七步：归一化相关的八倍误差。}
当 $\tau\le1$，参考矩阵为
\[
 G_0(L,J)=\frac13\left[R(K(L-J)/2)
 +\sum_{m\le M}b_m\sin(\pi mL)\sin(\pi mJ)\right],
 \quad R(t)=(1-|t|)_+^3.
\]
$R$ 是半正定核，因为
\[
 R(t)=\int_0^1 6(1-a)(a-|t|)_+da,
\]
且 $(a-|t|)_+=\langle\mathbf1_{[0,a]},\mathbf1_{[t,t+a]}\rangle$。
以上积分等式可在 $|t|\le1$ 时令 $u=1-|t|$，计算
$6\int_0^u(u-z)z\,dz=u^3$；$|t|>1$ 时两边均为0。
加上有限正谱仍半正定，且 $G_0(L,L)=F_M(L)/3\ge1/3$。
实际矩也是 Gram 矩阵，故其与参考矩阵的凸插值始终半正定，
并且对角至少 $1/4$。

对二乘二矩阵的归一化函数 $f(c,a,b)=c/\sqrt{ab}$，半正定性给
$|c|\le\sqrt{ab}$，故沿插值段
\[
 |\partial_cf|=1/\sqrt{ab}\le4,\qquad
 |\partial_af|=|c|/(2a\sqrt{ab})\le1/(2a)\le2,
 \qquad|\partial_bf|\le2.
\]
非对角误差至多 $\eta_N+e$；两个对角误差至多
$\eta_D+e\le\eta_N+e$。沿插值积分梯度，得
$|\kappa-\kappa_0|\le(4+2+2)(\eta_N+e)=8(\eta_N+e)$。

\textbf{第八步：短距和中距的相关缺口。}
令 $a_L=F_M(L)^{-1/2}$，
$v_L=(\sqrt{b_m}\sin(\pi mL)/\sqrt{F_M(L)})_{m\le M}$。
则 $a_L^2+\|v_L\|^2=1$，且
$\kappa_0=a_La_JR(\tau)+v_L\cdot v_J$。
单位向量 $(a_L,v_L)$ 和 $(a_J,v_J)$ 的 Cauchy 不等式给
\[
 1-\kappa_0=[1-a_La_J-v_L\cdot v_J]+a_La_J(1-R)
 \ge a_La_J(1-R).
\]
又 $F_M(L)\le\gamma L^2$，所以
$a_La_J\ge1/(\gamma LJ)\ge87/121$。
$0\le\tau\le1/8$ 时，
$1-R=\tau(3-3\tau+\tau^2)\ge2\tau$；
$1/8\le\tau\le1$ 时，$1-R\ge1-(7/8)^3=169/512$。
乘以 $87/121$ 并减去第七步误差，即为引理的两项缺口。
此外 $\|v_L\|^2=1-1/F_M(L)\le34/121$，$R\ge0$，故
$\kappa_0\ge-\|v_L\|\|v_J\|\ge-34/121$，短距、中距的实际相关
最终不小于 $-1/2$。

\textbf{第九步：远距的充分界和统一余量。}
当 $1\le\tau\le K/18$，差位移的 floor/next 均非零位移；
反射位移仍位于 $[K/9,2K/9]$，也具有此性质。
不需要把它们假设为满秩，第三步的异字符通用界已经覆盖全部缺陷：
\[
 |HH_{LJ}|\le4\mathcal C_{BK}=1/6+o(1).
\]
加上低谱界与完整 $e$ 误差，并用实际方差下界 $V_{\min}$，得到
\[
 \kappa\le\frac{4\mathcal C_{BK}+34/261+e}{V_{\min}},
 \qquad\kappa\ge-\frac{4\mathcal C_{BK}+1/116+e}{V_{\min}}.
\]
$B,K,M\to\infty$、$e\to0$ 时，所有误差对尺度一致，右端极限分别为
\[
 \frac{1/6+34/261}{100/261}=31/40,
 \qquad-\frac{1/6+1/116}{100/261}=-183/400.
\]
上界到 $4/5$ 的严格余量为 $1/40$，下界到 $-1/2$ 的余量为 $17/400$。
中距缺口的极限到 $1/5$ 的余量为
$14703/61952-1/5=11563/309760>0$。
因此可取同一个充分大阈值，使全安全集 $\kappa\ge-1/2$，
并且中、远距离缺口均至少 $1/5$。至此全部矩阵元、校正和归一化误差均已控制。
\end{proof}

\subsection{全部整数维数的实际密度}
\textbf{第一步：选择周期并验证全整数维数量词。}
先固定 $0<p<1$ 和有限 $w\ge4$；它们均不随 $n$ 增长。
置 $c_p=3p/\sqrt{29}$。对每个充分大整数 $n$，在 $K=2^d-1$ 中
选满足 $c_p^2n/K^2\ge16\log n$ 的最大者，令
\[
 v=\frac{c_p^2n}{K^2}=\frac{9p^2n}{29K^2},
 \qquad u=\sqrt v,\qquad h=v/2-2.
\]
此最大值存在：固定一个允许周期，在充分大 $n$ 时满足条件；而条件要求
$K^2\le c_p^2n/(16\log n)$，候选只有有限个。
所选周期的下一候选恰为 $2^{d+1}-1=2K+1$，它不满足条件。因此
\[
 \frac{c_p^2n}{(2K+1)^2}<16\log n\le\frac{c_p^2n}{K^2},
 \quad 16\log n\le v<(2+1/K)^2 16\log n.
\]
先仅用 $K\ge1$，可得 $v<144\log n$，从而
\[
 K=c_p\sqrt n/\sqrt v\ge\frac{c_p\sqrt n}{12\sqrt{\log n}}\to\infty.
\]
再次代回，上界可精化为 $v<(4+o(1))16\log n$。
特别地 $v=\Theta_p(\log n)$，而对每个固定 $r>0$，
\[
 \frac K{v^r}=\frac{c_p\sqrt n}{v^{r+1/2}}
 \ge\frac{c_p\sqrt n}{(144\log n)^{r+1/2}}\to\infty.
\]
这些不等式对每个充分大整数 $n$ 成立，并未限定周期切换的某个子序列。

\textbf{第二步：有限谱、标签分辨率和正宽峰。}
取不小于 $v$ 的最小二幂 $B$，$M=\lceil\sqrt v\rceil$，
$\delta=1/(wKv)$。则
\[
 v\le B<2v,\qquad M=O(\sqrt v),\qquad
 K+1>3B^4>H(B),\qquad1\le M<K
\]
最终同时成立，后两项由 $K/v^4\to\infty$ 推出。
因此每个大维数都存在引理\ref{lem:papersidon}的合格标签。
固定域表示、域元素排序和本原元后，可取合格有序基的字典序首项，
从而得到确定选择；存在性不要求有限随机尝试成功。

第 $j$ 个原子到所在格两端的距离均至少 $1/(2BK)$。
由于 $B<2v$、$w\ge4$，
\[
 \delta=\frac1{wKv}<\frac1{2BK},\qquad
 \frac\delta2<\frac1{4BK}<\frac1{2BK}.
\]
故整个箱峰留在本格内部，不同格的峰不重叠。箱平滑保持每个原子的总质量，
定义
\[
 \mu_{K,\delta}=K^{-1}\sum_{j=0}^{K-1}\beta_\delta(\cdot-a_j),
 \qquad \int\mu_{K,\delta}=1,
 \qquad0\le\mu_{K,\delta}\le\frac1{K\delta}=wv.
\]
箱宽对每个有限 $n$ 都严格为正，不需要先取原子极限。

\textbf{第三步：背景、低谱与实际密度的上下界。}
令 $c_m=\sqrt{2b_m/3}$、
\[
 D_M(x)=\sum_{m\le M}\pi m c_m\cos(2\pi mx),
 \qquad\rho(x)=1-p+p\mu_{K,\delta}(x)+\frac pK D_M(x).
\]
$D_M$ 的每个频率非零，故 $\int D_M=0$，于是 $\int\rho=1$。
系数界 $c_m\le\sqrt{32\alpha/3}/(\pi^2m^2)$ 给
\[
 \pi m c_m\le\frac{\sqrt{32\alpha/3}}{\pi m}<\frac{20}{3m},
 \qquad\|D_M\|_\infty\le\frac{20}3H_M,
 \quad H_M:=\sum_{m\le M}m^{-1}\le1+\log M.
\]
因 $H_M=O(\log v)$、$H_M/K\to0$，在固定 $p<1$ 后最终可使
$20pH_M/(3K)\le(1-p)/2$。由非负峰得到
\[
 \rho\ge1-p-\frac{20pH_M}{3K}\ge\frac{1-p}{2}>0,
\]
\[
 \rho\le1-p+pwv+\frac{20pH_M}{3K}=O_{p,w}(\log n).
\]
严格正下界依靠固定正背景 $1-p$。这里的高度界是单个一维因子的界，
并不声称其 $n$ 重乘积高度同阶。

\textbf{第四步：实际中心平均的准确公式。}
令 $\mathcal P_\delta=\mathcal P*\beta_\delta$。
由原函数分布导数和卷积可得
\[
 \mathcal P_\delta'=\sum_j\beta_\delta(\cdot-a_j)-K
 =K(\mu_{K,\delta}-1).
\]
所以对任意 $x,L$，周期区间积分满足
\[
 \int_{x-L/2}^{x+L/2}(\mu_{K,\delta}(y)-1)dy
 =\frac{\mathcal P_\delta(x+L/2)-\mathcal P_\delta(x-L/2)}K
 =\frac{H_{L,\delta}(x)}K.
\]
另一方面，直接积分余弦给
\[
 \frac1L\int_{x-L/2}^{x+L/2}\cos(2\pi my)dy
 =\frac{\sin(\pi mL)}{\pi mL}\cos(2\pi mx).
\]
平均密度中的 $\pi m$ 系数与分母准确相消，得到
\begin{equation}\label{eq:papersourceaverage}
 Y_L(x)=\frac1L\int_{x-L/2}^{x+L/2}\rho(y)dy=1+d_L(x),
 \qquad d_L=\frac p{KL}(H_{L,\delta}+B_L).
\end{equation}
平移差的积分为0，所有低频余弦积分也为0，故 $E_Pd_L=0$，$E_PY_L=1$。
全部尺度用的是同一个确定源 $\rho$。

\textbf{第五步：核预算与倾斜工具的统一资格。}
对这组参数，Gauss 交叉预算满足
\[
 C_x=O(K^{-1/2})+O((\log M)/K)=o(1/v),
 \qquad e=C_x+\frac2{3wv}.
\]
具体地，乘 $v$ 后第一项为 $O(v/\sqrt K)\to0$，第二项为
$O(v\log M/K)\to0$，均来自已证的 $K/v^r\to\infty$。
又 $T_M=O(v^{-3/2})$，
\[
 v\eta_D=O(v/B^2+v/K)\to0,
 \qquad v\eta_N=O(v/B^2+v/K)\to0.
\]
核引理的方差下界因此最终大于 $1/4$，并且
\[
 \limsup_n vr_*
 \le\frac{3\cdot2/(3w)}{100/87}=\frac{87}{50w}.
\]
令 $q_L=nE_Pd_L^2$。由准确平均公式，
\[
 q_L=\frac{np^2}{K^2L^2}G_\delta(L,L),\qquad
 v=\frac{np^2\gamma}{3K^2},\qquad
 \frac{q_L}{v}=\frac{3G_\delta(L,L)}{\gamma L^2}.
\]
所以所有安全尺度上同时有 $|q_L-v|\le vr_*=O_w(1)$。

箱卷积为正概率平滑，故 $|H_{L,\delta}|\le2$。
由 $c_m<20/(9m^2)$、$\sum_{m\ge1}m^{-2}=\pi^2/6<5/3$，
\[
 |B_L|\le\frac{20}9\sum_{m\le M}m^{-2}<\frac{100}{27},
 \qquad|H_{L,\delta}+B_L|<\frac{154}{27}.
\]
使用 $L\ge10/9$，得到统一实际幅度
\[
 |d_L|\le\frac pK\frac9{10}\frac{154}{27}
 =a_n:=\frac{77p}{15K}.
\]
若 $\sigma_L^2=E_Pd_L^2$，则
$\sigma_L=p\sqrt{G_\delta(L,L)}/(KL)$，因而
\[
 \frac{|d_L|}{\sigma_L}\le
 \frac{154/27}{\sqrt{G_\delta(L,L)}}\le\frac{308}{27}<14.
\]
最后，$K=c_p\sqrt n/\sqrt v$ 给
\[
 a_nv=O_p(v^{3/2}/\sqrt n)\to0,
 \qquad v^2/n=O_p((\log n)^2/n)\to0.
\]
全部阈值只依赖已固定的 $p,w$，对该维数的整个安全尺度集合一致。

% 两条标准外部概率定理须由主文件引言/参考文献正式注明来源。
\subsection{倾斜积分与亲和度}
以下采用 iid Berry--Esseen 定理\cite{Shevtsova2011}：若独立同分布实变量均值为 $\mu$、
方差为 $t^2>0$、第三绝对中心矩为 $m_3$，则标准化和的分布函数与
标准正态分布函数之差至多 $C_{\rm BE}m_3/(t^3\sqrt n)$，
其中可取 $C_{\rm BE}=0.4748$。
另采用独立向量的凸集正态比较\cite{Raic2019}：对总协方差为单位阵的独立中心化
$k$ 维向量 $V_a$，凸 Borel 集的比较误差至多
$(42k^{1/4}+16)\sum_aE\|V_a\|^3$。
我们在下述长引理中证明实际用到的奇异协方差扩展。

\begin{lemma}[加权尾的恒等式]\label{lem:paperweighted}
设 $Y>0$、$E_PY=1$，令 $dQ=YdP$、$S=\sum_{a=1}^n\log Y(X_a)$。
若 $Q^{\otimes n}$ 下 $S$ 的分布函数与 $N(M,s^2)$ 的分布函数相差至多
$\epsilon$，则任意实数 $h$ 有
\[
 e^hP^{\otimes n}(S>h)
 =\int_0^\infty e^{-r}Q^{\otimes n}(h<S\le h+r)dr,
\]
\[
 \left|e^hP^{\otimes n}(S>h)
 -\int_0^\infty e^{-x}\frac{\phi((h+x-M)/s)}s dx\right|\le2\epsilon.
\]
\end{lemma}
\begin{proof}
变测度使左端成为 $E_Q[e^{-(S-h)}\mathbf1_{\{S>h\}}]$。
对每个 $S>h$，$e^{-(S-h)}=\int_{S-h}^\infty e^{-r}dr$。
非负 Tonelli 交换积分即得第一式。区间概率是两个右连续分布函数之差，
两个端点各产生一次误差，$\int_0^\infty e^{-r}dr=1$。对正态区间概率再次
用 Tonelli，得到第二式中的密度积分。取严格不等式准确排除了 $S=h$ 的原子。

\medskip\noindent\textbf{积分端点与正态项的计算。}
为将上述等式写成逐步可核验的形式，记 $F_Q$ 为 $S$ 在 $Q^{\otimes n}$ 下的
右连续分布函数，$F_G(t)=\Phi((t-M)/s)$。则
\[
 Q(h<S\le h+r)=F_Q(h+r)-F_Q(h),
\]
故
\[
 \left|[F_Q(h+r)-F_Q(h)]-[F_G(h+r)-F_G(h)]\right|
 \le |F_Q(h+r)-F_G(h+r)|+|F_Q(h)-F_G(h)|\le2\epsilon.
\]
乘 $e^{-r}$ 积分时没有额外的长度因子。对正态部分，积分区域是
$0\le x\le r<\infty$，从而
\begin{align*}
 \int_0^\infty e^{-r}[F_G(h+r)-F_G(h)]\,dr
 &=\int_0^\infty\int_0^r e^{-r}\frac1s
       \phi\!\left(\frac{h+x-M}{s}\right)dx\,dr\\
 &=\int_0^\infty\left(\int_x^\infty e^{-r}dr\right)
       \frac1s\phi\!\left(\frac{h+x-M}{s}\right)dx.
\end{align*}
这给出所需表达式。所有被交换的函数非负，因此不预先假设 $S$ 有密度，
也不需要额外的绝对可积性论证。

\end{proof}

\begin{lemma}[保留亲和度的双点上界]\label{lem:paperpair}
设 $Y,Z>0$，$E_PY=E_PZ=1$，令
\[
 H=E_P\sqrt{YZ},\quad dQ=H^{-1}\sqrt{YZ}dP,
 \quad V=(\log Y+\log Z)/2,
\]
$t^2=\operatorname{Var}_QV>0$，$m_4=E_Q(V-E_QV)^4<\infty$，
$s=\sqrt n\,t$，$M=nE_QV$，$\epsilon=C_{\rm BE}\sqrt{m_4}/(t^2\sqrt n)$。
对同一组坐标的事件 $E_Y=\{\sum\log Y>h\}$、$E_Z=\{\sum\log Z>h\}$，
\[
 e^hP^{\otimes n}(E_Y\cap E_Z)
 \le H^n\left[\int_0^\infty e^{-r}\frac{\phi((h+r-M)/s)}s dr+2\epsilon\right]
 \le H^n[\phi(0)/s+2\epsilon].
\]
\end{lemma}
\begin{proof}
交事件包含于 $\{\sum V>h\}$，且变测度给
$e^hP(\sum V>h)=H^nE_Q[e^{-(\sum V-h)}\mathbf1_{\{\sum V>h\}}]$。
同引理\ref{lem:paperweighted}将后一个期望写成加权区间概率积分。
由 Cauchy，第三绝对中心矩至多 $t\sqrt{m_4}$，Berry--Esseen 误差至多
$\epsilon$。两个分布函数端点共计 $2\epsilon$，正态密度至多 $\phi(0)/s$。
全部比较误差始终位于实际亲和度 $H^n$ 的因子内。

\medskip\noindent\textbf{归一化与误差因子的保留。}
Cauchy 不等式给 $0<H\le\sqrt{E_PY\,E_PZ}=1$，所以 $Q$ 确为概率测度。
在乘积空间上，其密度恰为
\[
 \frac{dQ^{\otimes n}}{dP^{\otimes n}}=H^{-n}e^{\sum_aV(X_a)}.
\]
设 $R=\sum_aV(X_a)$，则
\begin{align*}
 e^hP(E_Y\cap E_Z)
 &\le e^hP(R>h)\\
 &=H^nE_{Q^{\otimes n}}[e^{-(R-h)}\mathbf1_{R>h}].
\end{align*}
由于
\[
 E_Q|V-E_QV|^3
 =E_Q\bigl(|V-E_QV|\,|V-E_QV|^2\bigr)
 \le\sqrt{E_Q|V-E_QV|^2}\sqrt{E_Q|V-E_QV|^4}
 =t\sqrt{m_4},
\]
Berry--Esseen 误差正是所列 $\epsilon$。须先在 $Q$ 下比较区间概率，
再乘回 $H^n$；若提前把此因子用1代替，随后对许多尺度对求和时便失去必要的衰减。

\end{proof}

在实际源上始终令 $X_1,\ldots,X_n$ 为独立 Haar 坐标，并定义
\[
 S_L=\sum_{a=1}^n\log Y_L(X_a),\qquad E_L=\{S_L>h\}.
\]
不同尺度事件共享同一组坐标。

\begin{proposition}[尺度对的统一概率估计]\label{prop:paperallpairs}
采用式\eqref{eq:papersourceaverage}的同一个源，固定 $p,w$ 后，全部安全
尺度对满足
\[
 ue^hP(E_L\cap E_J)\le C H_{LJ}^n,
 \quad H_{LJ}^n\le
 \exp\!\left[-\frac{(v-C)(1-\kappa(L,J))}{4(1+a_n)}\right].
\]
常数对尺度和 $n$ 一致。令 $c=174/121$。对每个固定 $r>0$，
\[
 \limsup_{n\to\infty}\sup_{|L-J|=2r/(Kv)}ue^hP(E_L\cap E_J)
 \le\phi(0)e^{4/(3w)}e^{-cr/4}.
\]
全部 近距 距离 $\tau=r/v\le1/8$ 另有可求和控制界 $Ce^{-cr/8}$。
中、远距离 $1/8\le\tau\le K/18$ 最终均有控制界 $Ce^{-v/25}$。
\end{proposition}
\begin{proof}
写 $Y=1+d$、$Z=1+e$，$T=(d+e)/2$。由 $\kappa\ge-1/2$ 和
$q_Y,q_Z=v+O(1)$，
\[
 s_0^2=E_PT^2=\frac{q_Y+q_Z+2\sqrt{q_Yq_Z}\kappa}{4n}
 \ge\frac{E_Pd^2+E_Pe^2}{8}.
\]
一致 Taylor 界给 $\|V-T\|_{L^2(P)}\le Ca_ns_0$。
由于 $dQ/dP=1+O(a_n)$，且 $E_PT=0$，
$\operatorname{Var}_QT=s_0^2(1+O(a_n))$。
中心化 $L^2$ 三角不等式遂给
\[
 n\operatorname{Var}_QV
 =\frac{q_Y+q_Z+2\sqrt{q_Yq_Z}\kappa}{4}(1+O(a_n)),
 \qquad s\asymp u.
\]
又 $|V-E_QV|\le2a_n/(1-a_n)$，因此
$m_4\le[2a_n/(1-a_n)]^2t^2$，$u\epsilon=O(a_n)\to0$。
成对引理给统一前因子。另一方面
\[
 1-H=\tfrac12E_P(\sqrt Y-\sqrt Z)^2
 \ge\frac{E_P(d-e)^2}{8(1+a_n)},
\]
而 $nE(d-e)^2=(\sqrt{q_Y}-\sqrt{q_Z})^2+
2\sqrt{q_Yq_Z}(1-\kappa)$。用 $H^n\le e^{-n(1-H)}$ 得所述指数上界。

固定 $r$ 时，由下一个局部增量引理，$nE(d_L-d_J)^2=O_{r,w}(1)$，
故 $\kappa\to1$ 一致；此时 $u/s\to1$，前因子精确到 $\phi(0)+o(1)$。
核引理的 近距 界给
\[
 v(1-\kappa)\ge cr-8v(\eta_N+e)=cr-\frac{16}{3w}-o(1).
\]
这证明固定 $r$ 的精确渐近式。
对全部增长 近距 距离，上一个下界仍一致，统一前因子结合
$(v-C)/(v(1+a_n))\ge1/2$ 给 $Ce^{-cr/8}$。
中远缺口至少 $1/5$，而上述比值最终至少 $4/5$，给 $Ce^{-v/25}$。
固定距离的精确渐近式与增长距离的控制界是两个不同断言，后者承担无穷求和。

\medskip\noindent\textbf{一致 Taylor 估计的细节。}
以下所有常数均不依赖尺度对。充分大 $n$ 时 $a_n\le1/2$。
由积分余项公式，对 $|z|\le a_n$ 有
\[
 |\log(1+z)-z|\le\frac{z^2}{2(1-a_n)},\qquad
 \left|\log(1+z)-z+\frac{z^2}{2}\right|
 \le\frac{|z|^3}{3(1-a_n)^3}.
\]
因此 $|V-T|\le C(d^2+e^2)\le Ca_n(|d|+|e|)$。
将其平方积分，并使用 $s_0^2\ge(E_Pd^2+E_Pe^2)/8$，得到
$\|V-T\|_2\le Ca_ns_0$。
又 $1-a_n\le\sqrt{YZ}\le1+a_n$、$1-a_n\le H\le1$，所以
$\|dQ/dP-1\|_\infty\le Ca_n$。
于是
\[
 E_QT^2=s_0^2+O(a_ns_0^2),\qquad
 |E_QT|=|E_P[(dQ/dP-1)T]|\le Ca_ns_0,
\]
从而 $\operatorname{Var}_QT=s_0^2(1+O(a_n))$。
对任意平方可积 $A,B$，中心化是正交投影，故
\[
 |\sqrt{\operatorname{Var}_QA}-\sqrt{\operatorname{Var}_QB}|
 \le\|(A-E_QA)-(B-E_QB)\|_{L^2(Q)}
 \le\|A-B\|_{L^2(Q)}.
\]
应用于 $A=V,B=T$ 即得到前面的方差比较。
此外，$V$ 的值域包含于
$[\log(1-a_n),\log(1+a_n)]$；值域宽度至多 $2a_n/(1-a_n)$，
中心化后仍受此宽度控制。因此
\[
 u\epsilon\le\frac{2C_{\rm BE}a_n}{1-a_n}\frac{u}{\sqrt n\,t}
 =O(a_n).
\]
最后，亲和度恒等式由 $EY=EZ=1$ 直接得到：
\[
 E(\sqrt Y-\sqrt Z)^2=2-2H,
 \qquad(\sqrt Y-\sqrt Z)^2
 =\frac{(d-e)^2}{(\sqrt Y+\sqrt Z)^2}
 \ge\frac{(d-e)^2}{4(1+a_n)}.
\]
这些计算说明前因子及指数界均对同一安全集上的所有尺度对一致。

\end{proof}

\subsection{原密度的局部增量与箱平滑核}
\begin{lemma}[原子位置差的计数]\label{lem:paperphysicalcount}
令 $N(\ell)$ 为圆周有序原子差恰为 $\ell/(BK)$ 的对数，写
$\ell=Bk+b_0$，$0\le b_0<B$。若 $k,k+1$ 均为非缺陷、非零位移 位移，则
\[
 N(\ell)=q/B-\mathbf1_{\{b_0=0\}}.
\]
此外 $N(0)=K$；若 $0<|\ell|<B$，则 $N(\ell)=q|\ell|/B^2$。
\end{lemma}
\begin{proof}
物理差来自 位移 $k$、标签差 $b_0$ 的 $B-b_0$ 个标签对，或 位移 $k+1$、
标签差 $b_0-B$ 的 $b_0$ 个标签对。用式\eqref{eq:paperlabelcount}相加得
$q/B$，仅 $b_0=0$ 的第一类含零标签对，须减1。
当 $0<\ell<B$，零位移 不能产生非零标签差，只有相邻 非缺陷 单元
的 $\ell$ 个标签对，给 $q\ell/B^2$。负差由交换有序对得到，零差恰为 $K$ 个 同原子。

\medskip\noindent\textbf{进位与零差的区分。}
用整数坐标 $z_j=Bj+y_j$ 表示原子位置（公共中点平移不影响差）。
若 $z_{j'}-z_j\equiv\ell\pmod{BK}$，则
\[
 B(j'-j)+(y_{j'}-y_j)=Bk+b_0\pmod{BK}.
\]
因 $-(B-1)\le y_{j'}-y_j\le B-1$，仅有两种可能：
$j'-j=k$、标签差为 $b_0$，或 $j'-j=k+1$、标签差为 $b_0-B$。
前一种有 $B-b_0$ 对标签，后一种有 $b_0$ 对标签；二者穷尽且互不相交。
当两个位移均为非缺陷非零位移时，每对标签的计数为 $q/B^2$，
只有 $(0,0)$ 标签对少一个，故
\[
 (B-b_0)\frac q{B^2}+b_0\frac q{B^2}-\mathbf1_{b_0=0}
 =\frac qB-\mathbf1_{b_0=0}.
\]
在 $0<\ell<B$ 的情形，第一种要求同一格内两个不同原子，但每格恰有一个原子，
所以该类计数为0；第二类的 $\ell$ 对标签均非零标签对，得到 $q\ell/B^2$。
当 $\ell=0$ 时，各原子只与自身重合，不能套用非零位移公式。

\end{proof}

\begin{lemma}[固定局部窗口的实际增量]\label{lem:paperlocal}
固定有限 $T$、$p\in(0,1)$、$w\ge4$。以安全基点 $L_0$ 置
$L(t)=L_0+2t/(Kv)$，$|t|\le T$，所用尺度均另要求安全。
令 $\varepsilon=1/w$，$\chi_\varepsilon$ 为 $[-\varepsilon/2,\varepsilon/2]$
上的均匀密度，$b_{[s,t]}=\mathbf1_{[s,t]}*\chi_\varepsilon$。
对固定有限节点 $t_j$，
\[
 D_{j\ell}:=nE_P(d_{L(t_j)}-d_{L(t_\ell)})^2
 =\lambda(L_0)\|b_{[0,t_j]}-b_{[0,t_\ell]}\|_2^2+o(1),
 \quad\lambda(L_0)=\frac{58}{9L_0^2}.
\]
全部余项对合格安全基点一致。基准 $d_{L_0}$ 与任一局部增量的
总协方差为 $O_T(1)$。
\end{lemma}
\begin{proof}
先冻结分母，置
\[
 U_{[s,t]}=\frac p{KL_0}(H_{L(t),\delta}-H_{L(s),\delta}).
\]
它由两个方向的平滑区间脉冲组成，每脉冲幅度系数为 $p/(KL_0)$，
物理长度单位为 $1/(Kv)$，并减去准确均值 $2p(t-s)/(K vL_0)$。
同方向 同原子 两族各有 $K$ 个脉冲，变量换元给系数
\[
 \frac{2np^2}{K^2vL_0^2}=\frac{58}{9L_0^2}.
\]
同方向异原子重合只涉及0附近宽 $O_T(1/(Kv))$ 的差窗口，故只有
$O_T(B/v+1)=O_T(1)$ 个整数 $\ell$，且最终 $|\ell|<B$。
前一引理给其总相对权重 $O_T(B^{-2})$。
反方向重合差落在 $L_0$ 附近同宽窗口。相应两 单元 位移 到
$K(L_0-1)$ 的距离不超过 $1+O_T(1/v)<2$，安全条件使二者均 非缺陷；
$I$ 又排除 零位移。每差格点相对权重至多 $q/(KB)$，全部镜像的总权重
为 $O_T(B^{-1})$。已包含跨 单元 进位及全部近邻镜像。
任意脉冲内积至多 $\sqrt{|[s,t]||[s',t']|}\le2T$，所以两种异原子误差均趋零。
均值乘积的 $n$ 倍为 $O_T(1/v)$。
因此
\[
 nE_PU_IU_J=\lambda(L_0)\langle b_I,b_J\rangle+o(1).
\]

低谱系数 $c_m=\sqrt{2b_m/3}$ 满足 $\sum m^2c_m^2<\infty$。
长度差 $O_T(1/(Kv))$，直接用 $|\sin x-\sin y|\le|x-y|$ 得低谱增量的
$n$ 倍平方误差 $O_T(1/(K^2v))$；将 $L(t)$ 分母换成 $L_0$ 也有同阶误差。
Minkowski 和 Cauchy 控制它们与 原子项 脉冲的交叉，不要求其正交。
极化即得 $D_{j\ell}$ 的公式。
原子项 增量的 $L^1$ 范数至多 $C_Tp/(Kv)$，基准幅度为 $O(p/K)$，
二者乘积再乘 $n$ 为 $O_T(1)$；低谱和分母项同样满足该界。

此外 $|\log(1+d_t)-\log(1+d_s)-(d_t-d_s)|\le Ca_n|d_t-d_s|$。
单倾斜或几何平均倾斜的密度均为 $1+O(a_n)$，所以替换为对数变量 的
倾斜协方差只产生 $o(1)$。均值平方项的 $n$ 倍至多 $O(v^2/n)$。
由 $a_nv\to0$，上述所有估计在固定窗口内一致。

\medskip\noindent\textbf{脉冲缩放与误差次数。}
设 $I_0=[s,t]$，$\ell_0=t-s\ge0$，令
$\psi_{I_0}(x)=b_{I_0}(Kvx)$。则
\[
 \int\psi_{I_0}\psi_{J_0}=\frac1{Kv}\langle b_{I_0},b_{J_0}\rangle,
 \qquad \int\psi_{I_0}=\frac{\ell_0}{Kv}.
\]
每个边界移动产生一族平移的 $\psi_{I_0}$；两端边界共两族。
因此同原子贡献乘 $n$ 后是
\[
 n\left(\frac p{KL_0}\right)^2\frac{2K}{Kv}
 \langle b_{I_0},b_{J_0}\rangle
 =\frac{2np^2}{K^2vL_0^2}\langle b_{I_0},b_{J_0}\rangle.
\]
一个非零差格点的项与一个同原子差格点的积分上界具有同一缩放，
其相对权重为 $N(\ell)/K$。
由于 $v\le B<2v$，固定窗口内可能的 $|\ell|$ 有统一的有限上界。
同向非零差满足
$N(\ell)/K=(q/K)|\ell|/B^2=O_T(B^{-2})$；
镜像差满足 $N(\ell)/K\le q/(KB)=O(B^{-1})$。
这直接给出所列两项误差，而没有以“几乎无重叠”代替计数。

下面核对低频与分母误差。Parseval 恒等式给
\[
 \|B_L-B_J\|_2^2
 \le\frac{\pi^2|L-J|^2}{2}\sum_{m\le M}m^2c_m^2
 \le C|L-J|^2.
\]
乘上 $np^2/K^2=29v/9$，再代入 $|L-J|=O_T((Kv)^{-1})$，
得到 $O_T((K^2v)^{-1})$。分母改动满足
$|L^{-1}-L_0^{-1}|\le C|L-L_0|$；因 $\|H_L+B_L\|_2=O(1)$，同样得到此界。
若把主脉冲写作 $U$，余项写作 $R$，则
\[
 |nE(U+R)^2-nEU^2|
 \le2\sqrt{nEU^2}\sqrt{nER^2}+nER^2=o(1).
\]
这也说明交叉项虽然不为零，仍可统一控制。

最后，由平移区间指示函数的 $L^1$ 变化量，
$\|H_{L(t),\delta}-H_{L(s),\delta}\|_1\le C_T/v$。
乘基准幅度与源系数，得到
\[
 n\left|E\left[d_{L_0}\frac p{KL_0}
       (H_{L(t),\delta}-H_{L(s),\delta})\right]\right|
 \le C_T\frac{np^2}{K^2v}=O_T(1).
\]
低频和分母余项用 Cauchy 控制后至多 $O_T(K^{-1})$。
这证明了条件高斯计算所需的基准与增量协方差有界性。

\end{proof}

\subsection{有限面比较与实际最大值主项}
\begin{lemma}[允许奇异协方差的有限面比较]\label{lem:paperpolyface}
设 $V_1,\ldots,V_n$ 独立同分布于 $\mathbb R^k$，均值为 $\mu$，
$\|V_1-\mu\|\le B_n$，$C=n\operatorname{Cov}(V_1)$。
令 $G\sim N(n\mu,C)$，允许 $C$ 奇异。设 $D$ 是 $J$ 个开或闭半空间的交，
单位法向为 $a_\ell$，每个 $b_\ell=a_\ell^TCa_\ell>0$。
对任意 $\sigma,\eta>0$，
\[
 |P(\sum_aV_a\in D)-P(G\in D)|
 \le4C_k\left(\frac{kB_n}{\sigma}+\frac{\gamma_k}{\sqrt n}\right)
 +4Je^{-\eta^2/(2\sigma^2)}
 +2\eta\sum_{\ell=1}^J(2\pi b_\ell)^{-1/2},
\]
$C_k=42k^{1/4}+16$，$\gamma_k=E\|Z_k\|^3$，$Z_k\sim N(0,I_k)$。
此上界与各面阈值无关。
\end{lemma}
\begin{proof}
每个中心化向量加独立 $\sigma Z_k/\sqrt n$，总噪声为 $\sigma Z_k$。
令 $A_\sigma=(C+\sigma^2I)^{-1/2}$ 白化，则
\[
 nE\|A_\sigma(V_1-\mu)\|^3
 \le\frac{B_n}{\sigma}nE\|A_\sigma(V_1-\mu)\|^2
 \le\frac{kB_n}{\sigma}.
\]
噪声部分的三阶矩和至多 $\gamma_k/\sqrt n$。用 $(a+b)^3\le4(a^3+b^3)$
与前述凸集比较，得第一项。
所有单位面法向的噪声投影绝对值超过 $\eta$ 的概率至多
$2Je^{-\eta^2/(2\sigma^2)}$。实际与加噪向量、加噪正态与原正态两次
包含比较共计入两次坏事件。各面阈值合计移动 $2\eta$，原正态面投影的
密度至多 $(2\pi b_\ell)^{-1/2}$，所以得最后一项。
可从扩张/收缩半空间的夹逼保留原来的开闭面；实际原子和平局不需零边界假设。

\medskip\noindent\textbf{半空间夹逼的具体形式。}
将 $D$ 写为 $a_\ell\cdot x\le c_\ell$ 的交（严格面按同样方式处理），
记 $D_t$ 为各阈值同时变成 $c_\ell+t$ 后的集合。
设 $S=\sum V_a$、$Z\sim N(0,I_k)$ 独立，
$\beta=2J\exp[-\eta^2/(2\sigma^2)]$。投影噪声均不超过 $\eta$ 的事件上，
\[
 S\in D\Longrightarrow S+\sigma Z\in D_\eta,
 \qquad S+\sigma Z\in D_{-\eta}\Longrightarrow S\in D.
\]
记加噪后的凸集比较误差为 $\Delta$。依次比较并去除正态噪声，得到
\[
 P(G\in D_{-2\eta})-\Delta-2\beta
 \le P(S\in D)
 \le P(G\in D_{2\eta})+\Delta+2\beta.
\]
$D_{2\eta}\setminus D$ 包含于各面外侧宽 $2\eta$ 的条带之并；
$D\setminus D_{-2\eta}$ 类似。每一侧的误差分别至多
$2\eta\sum_\ell(2\pi b_\ell)^{-1/2}$，取两侧较大值即可，
不需要再把两个条带界相加。

为核对 $\Delta$，令 $U_a=A_\sigma(V_a-\mu)$、
$Z_a'=A_\sigma\sigma Z_a/\sqrt n$。则
\[
 \sum_a\operatorname{Cov}(U_a+Z_a')=I_k,\quad
 \|A_\sigma\|\le\sigma^{-1},\quad
 \operatorname{tr}(A_\sigma C A_\sigma)\le k.
\]
从而 $\sum_aE\|U_a\|^3\le kB_n/\sigma$，
而 $\|A_\sigma\sigma\|\le1$ 给
$\sum_aE\|Z_a'\|^3\le\gamma_k/\sqrt n$。
凸集定理适用于白化后的集合；可逆线性映射保持凸性，故得所列 $\Delta$。
原协方差奇异不会妨碍 $C+\sigma^2I$ 的可逆性。

\end{proof}

\begin{theorem}[固定有限尺度段的越阈最大值等式]\label{thm:paperwinner}
固定 $0<p<1$、有限节点 $0=t_0<t_1<\cdots<t_{k-1}\le T$，$k\ge2$，
最小相邻距离 $A>0$，以及有限 $w\ge\max(4,1/A)$。
采用前述全部整数维数的源，设 $L_j=L_0+2t_j/(Kv)$ 全部在 $I$ 内且安全，
$S_j=\sum_{a=1}^n\log Y_{L_j}(X_a)$，$U=\{\max_jS_j>h\}$。
令 $b_t=\mathbf1_{[0,t]}*\chi_\varepsilon$，$\varepsilon=1/w$，$b_0=0$。
若 $W$ 是 $L^2(\mathbb R)$ 上的实等距 高斯 映射，定义
\[
 X_{\varepsilon,\lambda}(t)=\sqrt\lambda W(b_t)-\tfrac\lambda2\|b_t\|_2^2,
 \qquad\lambda(L_0)=\frac{58}{9L_0^2}.
\]
先固定上述全部参数，再令全部整数 $n\to\infty$，则
\begin{equation}\label{eq:paperwinner}
 ue^hP^{\otimes n}(U)=\phi(0)E\exp(\max_jX_{\varepsilon,\lambda(L_0)}(t_j))+o(1),
\end{equation}
余项对全部合格安全基点一致。
\par 随后保持有限节点，令 $w\to\infty$，右侧期望在
$\lambda\in[522/121,522/100]$ 上一致收敛到
\[ E\exp(\max_j[\sqrt\lambda B(t_j)-\lambda t_j/2]).\]
\end{theorem}
\begin{proof}
将平局分给最小索引：$j$ 胜出当且仅当
$S_\ell<S_j$ 对 $\ell<j$，$S_\ell\le S_j$ 对 $\ell>j$。
对每个实际向量，这些域构成不交分解。
每个胜出项使用自己的单倾斜 $dQ_j=Y_{L_j}dP$。
变测度和非负 Tonelli 给
\begin{equation}\label{eq:paperwinneridentity}
 e^hP(U)=\sum_j\int_0^\infty e^{-r}
 Q_j^{\otimes n}\{h<S_j\le h+r,\ j\text{胜出}\}dr.
\end{equation}
积分域是 $k+1$ 个开闭半空间的交。

写 $d_j=Y_{L_j}-1$，$g_j=\log Y_{L_j}$，
$q_j=nEd_j^2$，$C^P_{j\ell}=nEd_jd_\ell$，
$D_{j\ell}=nE(d_j-d_\ell)^2$。
局部引理给 $D_{j\ell}=\lambda\|b_{t_j}-b_{t_\ell}\|_2^2+o(1)$。
直接在实际 $Q_j$ 下作二阶展开，所有误差为 $O(a_nv+v^2/n)=o(1)$，得到
\begin{align*}
 nE_{Q_j}g_\ell&=C^P_{j\ell}-q_\ell/2+o(1),&nE_{Q_j}g_j&=q_j/2+o(1),\\
 nE_{Q_j}(g_\ell-g_j)&=-D_{j\ell}/2+o(1),\\
 n\operatorname{Cov}_{Q_j}(g_\ell-g_j,g_m-g_j)
 &=(D_{j\ell}+D_{jm}-D_{\ell m})/2+o(1),\\
 n\operatorname{Cov}_{Q_j}(g_j,g_\ell-g_j)
 &=(q_\ell-q_j-D_{j\ell})/2+o(1)=O(1).
\end{align*}
例如均值来自 $E[(1+d_j)(d_\ell-d_\ell^2/2)]$，
三阶余项的 $n$ 倍至多 $Ca_nv$；协方差的均值乘积误差至多 $Cv^2/n$。
差均值中 $q_\ell-q_j$ 相消，不需设所有实际方差等于 $v$。
基准总方差为 $q_j+o(1)=v+O(1)$。

若 $s=|t_j-t_\ell|\ge A\ge\varepsilon$，分段线性箱积分给
\[
 \|b_{t_j}-b_{t_\ell}\|_2^2=s-\varepsilon/3\ge2A/3.
\]
这里的箱积分可以直接分三段验证。平移后取区间 $[0,s]$，$s\ge\varepsilon$，
其与宽 $\varepsilon$ 的均匀箱卷积在左右两个长度 $\varepsilon$ 的边缘上线性变化，
在中间长度 $s-\varepsilon$ 的区间上等于1。因此
\[
 \|\mathbf1_{[0,s]}*\chi_\varepsilon\|_2^2
 =2\int_0^\varepsilon(x/\varepsilon)^2dx+(s-\varepsilon)
 =s-\varepsilon/3.
\]
又 $b_{t_j}-b_{t_\ell}$ 正是区间 $[t_\ell,t_j]$ 的卷积（或其相反数），
故平方范数只依赖两节点间距。

因此比较面 $(e_\ell-e_j)/\sqrt2$ 的方差最终至少
$\lambda_{\min}A/6$，两基准面的方差最终至少 $v/2$。
每个中心化单坐标 log 向量的范数至多 $2\sqrt k\,a_n/(1-a_n)$。
取 $\sigma=n^{-1/8}$、$\eta=\sigma\sqrt{8\log n}$，有限面比较的误差乘 $u$ 为
\[
 O_{k,A,T,w,p}(vn^{-3/8}+u/\sqrt n+un^{-4}
 +un^{-1/8}\sqrt{\log n})=o(1).
\]
误差与阈值无关，故可逐 $r$ 积分；有限 $k$ 项相加仍为 $o(1/u)$。
这一步只要求各面非退化，允许整个差 Gram 奇异。

对第 $j$ 个倾斜测度的矩定义正态向量 $G^{(j)}$，仅条件化其基准坐标
$G_j^{(j)}=h+x$。基准方差 $\sim v$，$h-q_j/2=O(1)$，基准与差向量
协方差为 $O(1)$。对固定 $x\ge0$，条件均值与协方差的修正遂趋零。
只除以基准的一维方差，无需逆整个差 Gram。
极限差向量恰是有限 高斯 向量 $X_{\varepsilon,\lambda}$
按 $e^{X_{\varepsilon,\lambda}(t_j)}$ 倾斜后的差向量：均值为
$-\lambda\|b_{t_\ell}-b_{t_j}\|_2^2/2$，协方差是上述极化主项。
此等式由 高斯 指数倾斜公式直接计算而来。
每个不同节点之差有正方差，故极限胜出域的边界是有限个零概率超平面，
即使差向量奇异，弱收敛仍给该域概率的收敛。

令 $f_{j,n}$ 为基准正态密度，$\pi_{j,n}(x)$ 为条件胜出概率。
对薄层再次用 Tonelli，得到其加权积分为
\[
 \Gamma_{j,n}=\int_0^\infty e^{-x}f_{j,n}(h+x)\pi_{j,n}(x)dx.
\]
$uf_{j,n}(h+x)$ 一致有界，并对固定 $x$ 趋于 $\phi(0)$；$e^{-x}$ 是可积控制界。
因此 $u\Gamma_{j,n}\to\phi(0)P_j(j\text{胜出})$。
所有实际矩余项对基点一致，$\lambda$ 的取值区间紧。
若一致收敛失败，则从失败基点子列抽取 $\lambda$ 收敛子列；以上矩收敛、
面非退化及参数连续性再次给同一极限，矛盾。这证明安全基点的一致性。
最后，有限 高斯 的准确倾斜与分区给
\[
 \sum_jP_j(j\text{胜出})
 =\sum_jE[e^{X_j}\mathbf1_{\{j\text{胜出}\}}]=E e^{\max_jX_j},
\]
从而得式\eqref{eq:paperwinner}。

保持有限节点时，$b_t\to\mathbf1_{[0,t]}$ 于 $L^2$，有限向量收敛到
所述 Brownian 向量，且
\[
 E e^{2\max_jX_j}\le\sum_jE e^{2X_j}
 =\sum_je^{\lambda\|b_{t_j}\|_2^2}\le k e^{\lambda_{\max}T}.
\]
这控制指数最大值的一致可积性；同一紧参数子列论证给期望的一致收敛。

\medskip\noindent\textbf{条件化与指数倾斜的矩计算。}
为明确高斯极限的识别，写
$K_{j\ell}=\lambda\langle b_{t_j},b_{t_\ell}\rangle$。
原高斯向量 $X$ 的均值为 $-K_{jj}/2$，协方差为 $K$。
因 $Ee^{X_j}=1$，$e^{X_j}dP$ 是概率测度；在此测度下，均值变为
$-K_{\ell\ell}/2+K_{j\ell}$，协方差仍为 $K$。于是
\begin{align*}
 E_j(X_\ell-X_j)
 &=-\tfrac12(K_{\ell\ell}+K_{jj}-2K_{j\ell})=-D_{j\ell}^{\infty}/2,\\
 \operatorname{Cov}_j(X_\ell-X_j,X_m-X_j)
 &=K_{\ell m}-K_{\ell j}-K_{jm}+K_{jj}\\
 &=\tfrac12(D_{j\ell}^{\infty}+D_{jm}^{\infty}-D_{\ell m}^{\infty}).
\end{align*}
其中 $D_{j\ell}^{\infty}=\lambda\|b_{t_j}-b_{t_\ell}\|_2^2$。
这些正是实际倾斜矩展开给出的极限。

对实际正态向量，令基准均值、方差为 $m_j,s_j^2$，差向量均值为 $m_\Delta$，
基准与差向量的协方差列为 $c_j$。给定基准值 $h+x$ 后，条件矩为
\[
 m_\Delta+\frac{c_j}{s_j^2}(h+x-m_j),\qquad
 C_\Delta-\frac{c_jc_j^T}{s_j^2}.
\]
此处 $c_j=O(1)$、$s_j^2=v+O(1)$、$h-m_j=O(1)$。
所以对固定 $x$，两项修正均为 $O_x(v^{-1})$。
同时
\[
 uf_{j,n}(h+x)=\frac{u}{s_j}\phi\!\left(\frac{h+x-m_j}{s_j}\right)
 \longrightarrow\phi(0),\qquad
 uf_{j,n}(h+x)\le C.
\]
在 $x\in[0,R]$ 上可使用同一界，而 $x>R$ 的积分统一至多 $Ce^{-R}$。
先取 $n$ 极限、后令 $R\to\infty$，即可严谨地交换薄层积分与极限。

至于固定节点下令 $w\to\infty$，箱平均是 $L^2$ 近似恒等算子，
故 $\|b_t-\mathbf1_{[0,t]}\|_2\to0$。
在同一个等距高斯映射下，
$E|W(b_t)-W(\mathbf1_{[0,t]})|^2=\|b_t-\mathbf1_{[0,t]}\|_2^2$。
有限个节点因而联合依概率收敛；平方指数矩上界使指数最大值一致可积，
将依概率收敛提升为期望收敛。这一步不能在节点数未经控制地增长时直接使用。

\end{proof}

% 03片段的详细推导版本；保留原标签，无preamble，不覆盖原片段。
\subsection{Brownian 指数最大值及统一网格误差}
\begin{lemma}[带漂移最大值的指数均值]\label{lem:paperbrownian}
对标准 Brownian 运动 $B$，令 $X(t)=\sqrt2B(t)-t$，
$M_T=\sup_{0\le t\le T}X(t)$。对 $T>0$ 有
\[
 P(M_T>x)=\Phi\!\left(-\frac{x+T}{\sqrt{2T}}\right)
 +e^{-x}\Phi\!\left(\frac{T-x}{\sqrt{2T}}\right)\qquad(x\ge0),
\]
\[
 F(T):=Ee^{M_T}=(T+2)\Phi(\sqrt{T/2})
 +\sqrt{2T}\phi(\sqrt{T/2}),\qquad F(T)/T\longrightarrow1.
\]
\end{lemma}
\begin{proof}
先在零漂移测度 $P_0$ 下考虑坐标过程 $Z(t)=\sqrt2B(t)$。
其终点密度是
\[
 p_T(z)=\frac1{\sqrt{4\pi T}}\exp\!\left(-\frac{z^2}{4T}\right).
\]
因为 $Z(T)$ 为方差 $2T$ 的中心正态变量，
\[
 E_0\exp[-Z(T)/2-T/4]
 =\exp[(1/2)^2(2T)/2-T/4]=1.
\]
故可在 $\mathcal F_T$ 上定义概率测度
\[
 \frac{dP_{-}}{dP_0}=\exp[-Z(T)/2-T/4].
\]
对任意有限分割，将这个指数密度按独立增量分解。
\par 每段增量的零均值正态
密度乘以 $\exp[-\Delta Z/2-\Delta t/4]$ 后，完成平方即成为均值
$-\Delta t$、方差 $2\Delta t$ 的正态密度。因而在 $P_-$ 下，坐标过程
$Z$ 与原测度下的 $X(t)=\sqrt2B(t)-t$ 同分布；这里变测度的密度依赖
\emph{零漂移坐标过程的终点}，而非先把带漂移终点代入同一密度。

零漂移反射原理给出：对 $x>0$，在 $z<x$ 的终点纤维上，事件
$\{\sup_{t\le T}Z(t)>x\}$ 的联合终点密度为 $p_T(2x-z)$；
在 $z\ge x$ 的纤维上，其终点本身已达到该高度。因此
\begin{align*}
 P(M_T>x)
 &=\int_x^\infty e^{-z/2-T/4}p_T(z)\,dz
 +\int_{-\infty}^x e^{-z/2-T/4}p_T(2x-z)\,dz.
\end{align*}
第一项的指数满足
\[
 -\frac{z^2}{4T}-\frac z2-\frac T4
 =-\frac{(z+T)^2}{4T},
\]
所以它等于 $\Phi(-(x+T)/\sqrt{2T})$。
第二项中有
\[
 -\frac{(2x-z)^2}{4T}-\frac z2-\frac T4
 =-x-\frac{(z-2x+T)^2}{4T}.
\]
积分到 $z=x$ 后，得到 $e^{-x}\Phi((T-x)/\sqrt{2T})$。
这证明 $x>0$ 的反射公式。其右端在 $x\downarrow0$ 时趋于1，
故由事件的单调极限也得到 $P(M_T>0)=1$，公式在 $x=0$ 同样成立。

由于路径在 $t=0$ 的值为0，有 $M_T\ge0$。逐点写成
\[
 e^{M_T}=1+\int_0^\infty e^x\mathbf1_{\{M_T>x\}}\,dx,
\]
非负 Tonelli 定理给
\[
 F(T)=1+I_1(T)+I_2(T),
\]
其中，记 $z_T=\sqrt{2T}$、$a_T=T/z_T=\sqrt{T/2}$，
\[
 I_1(T)=\int_0^\infty e^x\Phi(-(x+T)/z_T)\,dx,
 \qquad I_2(T)=\int_0^\infty\Phi((T-x)/z_T)\,dx.
\]
先计算第一项。正态尾的 高斯 衰减保证
$e^x\Phi(-(x+T)/z_T)\to0$，所以分部积分没有无穷端边界项，且
\begin{align*}
 I_1(T)
 &=-\Phi(-a_T)+\frac1{z_T}\int_0^\infty
 e^x\phi((x+T)/z_T)\,dx\\
 &=-\Phi(-a_T)+\frac1{z_T}\int_0^\infty
 \phi((x-T)/z_T)\,dx\\
 &=-\Phi(-a_T)+\Phi(a_T)=2\Phi(a_T)-1.
\end{align*}
中间的完成平方使用
$x-(x+T)^2/(4T)=-(x-T)^2/(4T)$。
对第二项令 $y=(T-x)/z_T$。由于
$(y\Phi(y)+\phi(y))'=\Phi(y)$，且该原函数在 $y\to-\infty$ 时趋于0，
\begin{align*}
 I_2(T)&=z_T\int_{-\infty}^{a_T}\Phi(y)\,dy\\
 &=z_T\{a_T\Phi(a_T)+\phi(a_T)\}
 =T\Phi(a_T)+\sqrt{2T}\phi(a_T).
\end{align*}
合并两项即得所述 $F(T)$。最后 $a_T\to\infty$，
$\Phi(a_T)\to1$，$\sqrt{2/T}\,\phi(a_T)\to0$，因此 $F(T)/T\to1$。
\end{proof}

\begin{lemma}[指数加权的网格逼近]\label{lem:papermesh}
令 $\kappa\in[\kappa_-,\kappa_+]=[261/121,261/100]$，
$X_\kappa(t)=\sqrt{2\kappa}B(t)-\kappa t$。整数 $T\ge2$ 时取
$N=T^3$、$A=T^{-2}$，定义
\[
 H_\kappa(T,A)=E\exp\!\left(\max_{0\le j\le N}X_\kappa(jA)\right).
\]
则
\[
 \sup_{\kappa\in[\kappa_-,\kappa_+]}
 \left|\frac{H_\kappa(T,T^{-2})}{T}-\kappa\right|\longrightarrow0.
\]
\end{lemma}
\begin{proof}
Brownian 时间缩放使 $\sqrt{\kappa}B(t)$ 与 $B(\kappa t)$ 同分布，故连续
最大值 $M_c=\sup_{t\le T}X_\kappa(t)$ 的指数均值恰为
\[
 F_\kappa(T)=Ee^{M_c}=F(\kappa T).
\]
记网格最大值为 $M_g=\max_{0\le j\le N}X_\kappa(jA)$。
对每个节点 $t_j=jA$，$Ee^{X_\kappa(t_j)}=1$，因而
\[
 dQ_j=e^{X_\kappa(t_j)}dP,
 \qquad Q=\frac1{N+1}\sum_{j=0}^NQ_j,
 \qquad D_Q=\frac{dQ}{dP}
 =\frac1{N+1}\sum_{j=0}^Ne^{X_\kappa(t_j)}
\]
都是概率测度或概率密度。该混合的作用是控制指数最大值在坏路径上的
贡献；仅在原测度下估计坏路径的概率，并不足以控制其指数权重。

为确定混合下的路径分布，在 $Q_j$ 下分解独立 Brownian 增量。
节点 $t_j$ 之前的增量乘以指数倾斜
\(\exp[\sqrt{2\kappa}\,\Delta B-\kappa\Delta t]\)，
所以 $\Delta B$ 的新均值为 $\sqrt{2\kappa}\Delta t$，方差保持为
$\Delta t$。节点之后的增量不受倾斜。于是 $X_\kappa$ 在节点前后的
漂移分别为 $+\kappa$ 和 $-\kappa$，随机部分始终为方差速率
$2\kappa$ 的 Brownian 运动。由于锚点本身属于网格，每个网格段中
漂移是常数，其绝对值不超过 $\kappa_+$。

定义所有网格段相对左端的最大绝对变动
\[
 \Omega=\max_{0\le i<N}\sup_{0\le s\le A}
 |X_\kappa(iA+s)-X_\kappa(iA)|.
\]
置 $a=\kappa_+A$。在任一混合成分 $Q_j$ 下，若某一段的绝对变动
超过 $x>a$，则其无漂移 Brownian 部分的绝对最大值超过 $x-a$。
两侧反射原理、正态尾界以及全部 $N$ 段的并界给
\[
 Q_j(\Omega>x)\le4N\exp\!\left(-\frac{(x-a)^2}{4a}\right).
\]
同一个上界对所有 $j$ 成立，故均匀混合后仍有
\begin{equation}\label{eq:paperomegatail}
 Q(\Omega>x)\le4N\exp\!\left(-\frac{(x-a)^2}{4a}\right)
 \qquad(x>a).
\end{equation}

取 $\delta_T=8\sqrt{a\log T}$。对 $T\ge2$，
$\log T\ge\log2>2/3>\kappa_+/4\ge\kappa_+/T^2=a$，所以
$\delta_T\ge8a$。当 $x\ge\delta_T$ 时，
$x\ge(2+\sqrt2)a$，从而 $(x-a)^2\ge x^2/2$；式
\eqref{eq:paperomegatail}可放宽为
\[
 Q(\Omega>x)\le4Ne^{-x^2/(8a)}.
\]
对非负随机变量的指数尾积分恒等式给
\begin{align*}
 E_Q[e^\Omega\mathbf1_{\{\Omega>\delta_T\}}]
 &=e^{\delta_T}Q(\Omega>\delta_T)
 +\int_{\delta_T}^\infty e^xQ(\Omega>x)\,dx\\
 &\le4N e^{\delta_T-\delta_T^2/(8a)}
 +4N\int_{\delta_T}^\infty e^{x-x^2/(8a)}\,dx.
\end{align*}
函数 $g(x)=x-x^2/(8a)$ 的导数为 $1-x/(4a)$。由于
$\delta_T\ge8a$，对 $x\ge\delta_T$ 有
\[
 g'(x)\le1-\delta_T/(4a)\le-\delta_T/(8a).
\]
因此
\[
 \int_{\delta_T}^\infty e^{g(x)}\,dx
 \le e^{g(\delta_T)}\int_{\delta_T}^\infty
 e^{-\delta_T(x-\delta_T)/(8a)}\,dx
 =\frac{8a}{\delta_T}e^{g(\delta_T)}.
\]
再用 $\delta_T^2/(8a)=8\log T$，得到明确上界
\begin{equation}\label{eq:paperweightedomega}
 E_Q[e^\Omega\mathbf1_{\{\Omega>\delta_T\}}]
 \le4Ne^{\delta_T}T^{-8}\left(1+\frac{8a}{\delta_T}\right).
\end{equation}

逐路径有 $M_g\le M_c\le M_g+\Omega$，并且
$D_Q\ge e^{M_g}/(N+1)$。在好集 $\{\Omega\le\delta_T\}$ 上，
$e^{M_g}\ge e^{-\delta_T}e^{M_c}$，故
\[
 (e^{M_c}-e^{M_g})\mathbf1_{\{\Omega\le\delta_T\}}
 \le(1-e^{-\delta_T})e^{M_c}.
\]
在坏集上舍去非负的 $e^{M_g}$，再变到 $Q$：
\begin{align*}
 E[(e^{M_c}-e^{M_g})\mathbf1_{\{\Omega>\delta_T\}}]
 &\le E_Q\left[\frac{e^{M_c}}{D_Q}
 \mathbf1_{\{\Omega>\delta_T\}}\right]\\
 &\le(N+1)E_Q[e^\Omega\mathbf1_{\{\Omega>\delta_T\}}].
\end{align*}
结合式\eqref{eq:paperweightedomega}得
\begin{equation}\label{eq:papergriderror}
 0\le F_\kappa(T)-H_\kappa(T,A)
 \le(1-e^{-\delta_T})F_\kappa(T)
 +4N(N+1)e^{\delta_T}T^{-8}
 \left(1+\frac{8\kappa_+A}{\delta_T}\right).
\end{equation}

现 $\delta_T=O(\sqrt{\log T}/T)$，$e^{\delta_T}$ 有界，
$8\kappa_+A/\delta_T=O(1/(T\sqrt{\log T}))$，而
$N(N+1)T^{-8}=T^{-2}+T^{-5}$。因此第二项是 $O(T^{-2})$，
除以 $T$ 后趋于0。第一项除以 $T$ 后至多为
$\delta_T F(\kappa T)/T$，同样一致趋于0。
前一引理的准确公式又给
\[
 \frac{F(\kappa T)}T
 =\left(\kappa+\frac2T\right)\Phi(\sqrt{\kappa T/2})
 +\sqrt{\frac{2\kappa}{T}}\phi(\sqrt{\kappa T/2})
 \longrightarrow\kappa
\]
在 $\kappa\in[\kappa_-,\kappa_+]$ 上一致：正态尾的自变量至少为
$\sqrt{\kappa_-T/2}\to\infty$，其余系数在紧区间上一致有界。
这与式\eqref{eq:papergriderror}合起来证明所述网格极限。
\end{proof}

\section{主定理的证明}\label{sec:mainproof}
\subsection{安全尺度段的排列与交集估计}
固定 $0<p<1$、有限 $T=NA>0$、$D>0$ 和
$w\ge\max(4,1/A)$，其中 $N$ 为正整数。所有这些参数先固定，
再沿前述同一个原密度取整数维数极限。记 $P_0=T+D$，
原始基点及段内节点为
\[
 L_i=10/9+\frac{2iP_0}{Kv},\qquad
 L_{i,a}=L_i+\frac{2aA}{Kv},\quad0\le a\le N.
\]
只考虑末节点仍在 $I$ 内的 $i\ge0$；若任一节点不安全，则删除该整段。
保留的原始索引集合记为 $\mathcal I_n$，其势为 $J_*$。
对 $i\in\mathcal I_n$，记
\[
 U_i=\{\max_{0\le a\le N}S_{L_{i,a}}>h\}.
\]

\begin{lemma}[安全删除与所有原始距离的控制]\label{lem:paperpackets}
上述固定参数下，
\[
 J_*\sim\frac{Kv|I|}{2P_0},\qquad
 \frac{J_*}{u\sqrt n}\longrightarrow\frac{c_p|I|}{2P_0}.
\]
令 $c=174/121$。全部无序段对满足
\[
 \limsup_{n\to\infty}\frac{ue^h}{J_*}
 \sum_{\substack{i,j\in\mathcal I_n\\i<j}}P(U_i\cap U_j)
 \le\phi(0)e^{4/(3w)}(N+1)^2
 \frac{e^{-cD/4}}{1-e^{-cP_0/4}}.
\]
\end{lemma}
\begin{proof}
先准确计数未删除的原段。条件 $L_{i,N}\le11/9$ 等价于
$iP_0+T\le Kv|I|/2$，所以充分大 $n$ 时原段数为
\[
 J_0=\left\lfloor\frac{Kv|I|/2-T}{P_0}\right\rfloor+1
 =\frac{Kv|I|}{2P_0}+O_{T,D}(1).
\]
这里末节点造成的端点损失为有限的 $O(T/P_0+1)$，并非随 $n$ 增长的误差。

安全条件在 $x=K(L-1)$ 坐标中删去每个缺陷周围半径2的带。
由于 $L\in I$，其物理范围为 $x\in[K/9,2K/9]$，区间长度为 $K/9$。
当 $K$ 充分大时，每个周期缺陷在这一范围内至多有一个相关副本；
若一个删除带与区间端点相交，仍按同一个带计数。
映回长度坐标，带宽至多为 $4/K$。
固定段内索引 $a$ 后，$L_{i,a}$ 随 $i$ 构成间距
$\Delta_n=2P_0/(Kv)$ 的等距网格。任一宽 $4/K$ 的带中至多有
\[
 \frac{4/K}{\Delta_n}+2=\frac{2v}{P_0}+2
\]
个这样的节点。对全部 $N+1$ 个段内节点和全部缺陷作并计数，删除段数
不超过
\[
 (N+1)|\mathcal D|\left(\frac{2v}{P_0}+2\right).
\]
最小二幂标签尺度满足 $v\le B<2v$，而
$|\mathcal D|=(B-1)(B-2)<4v^2$。故删除数是
$O_{T,A,D}(v^3)$；重叠带或一个段含多个坏节点只会使此上界变大。
增长处方已给 $K/v^r\to\infty$ 对每个固定 $r>0$ 成立，特别是
$v^3/(Kv)=v^2/K\to0$。因此
\[
 J_*=J_0-O_{T,A,D}(v^3)
 \sim\frac{Kv|I|}{2P_0}.
\]
再用 $v=u^2$ 和准确关系 $K=c_p\sqrt n/u$，即得
$J_* /(u\sqrt n)\to c_p|I|/(2P_0)$。

下面按\emph{原始索引差}分类全部交集。若保留段的原索引为 $i<j$，
写 $q=j-i\ge1$，则对其任意节点 $a,b$，
\begin{align*}
 L_{j,b}-L_{i,a}
 &=\frac{2\{qP_0+(b-a)A\}}{Kv},\\
 r_{q,a,b}:=qP_0+(b-a)A
 &\ge qP_0-T=D+(q-1)P_0>0.
\end{align*}
段的先后次序因 $D>0$ 而在全部节点之间保持一致。
此等式不要求 $P_0/A$ 为整数。删段后的相邻保留序号可能对应很大的
原始网格距离 $q$；将其替换为新序号差会改变实际距离，故始终保留原索引。
对每个固定 $q$，每个保留 $i$ 至多对应一个 $j=i+q$，所以段对数至多为 $J_*$。
又
\[
 U_i\cap U_j\subseteq
 \bigcup_{a,b=0}^N(E_{L_{i,a}}\cap E_{L_{j,b}}),
\]
故必须保留全部 $(N+1)^2$ 个节点交集，不能只选段首或一个代表节点。

写 $\tau=r_{q,a,b}/v$。先处理近区 $0<\tau\le1/8$。
对每个固定 $q$ 及固定 $a,b$，$r_{q,a,b}$ 是有限正数，因 $v\to\infty$，
该节点对最终属于近区。命题\ref{prop:paperallpairs}的固定距离精确渐近界给
\[
 \limsup_n\sup_{
 \substack{L,J\text{安全}\\J-L=2r_{q,a,b}/(Kv)}}
 ue^hP(E_L\cap E_J)
 \le\phi(0)e^{4/(3w)}e^{-cr_{q,a,b}/4}.
\]
若某维数无这样的实际节点对，该项直接记为0，不从空集合取有利下确界。
于是先固定整数 $Q$，有限求和与每个原始网格距离 $q$ 的段对数上界给
\begin{align*}
 \limsup_n\frac{ue^h}{J_*}
 \sum_{\substack{i<j\\j-i\le Q}}P(U_i\cap U_j)
 &\le\phi(0)e^{4/(3w)}
 \sum_{q=1}^Q\sum_{a,b=0}^Ne^{-cr_{q,a,b}/4}\\
 &\le\phi(0)e^{4/(3w)}(N+1)^2
 \sum_{q=1}^Qe^{-c[D+(q-1)P_0]/4}.
\end{align*}
有限 $Q$ 的全部相关节点对最终都在近区，故此步可以直接使用锋利极限。

对 $q>Q$ 而仍属于近区的节点对，使用同命题的\emph{有限维统一控制界}
$Ce^{-cr/8}$。常数可依赖全部已固定参数，却不依赖 $n,q,L,J$。
因此其归一化总误差不超过
\[
 C(N+1)^2\sum_{q>Q}e^{-c[D+(q-1)P_0]/8}.
\]
该几何尾和随 $Q\to\infty$ 趋于0。先取有限 $Q$ 下的 $n$ 极限，
再令 $Q\to\infty$，得到整个近区的上界
\[
 \phi(0)e^{4/(3w)}(N+1)^2
 \sum_{q\ge1}e^{-c[D+(q-1)P_0]/4}
 =\phi(0)e^{4/(3w)}(N+1)^2
 \frac{e^{-cD/4}}{1-e^{-cP_0/4}}.
\]
这一步将固定距离的极限与增长距离的控制界分开；没有把第一个式子的
$o(1)$ 当作对所有增长 $q$ 一致的误差。

对每个固定 $a,b$，中区 $v/8\le qP_0+(b-a)A\le v$ 对应的原始网格距离 $q$
数至多为 $7v/(8P_0)+2=O_{T,A,D}(v)$。每项有
$ue^hP(E_L\cap E_J)\le Ce^{-v/25}$，故中区归一误差至多
$O_{p,w,T,A,D}(v)e^{-v/25}\to0$。
全部长度在 $I$ 内又保证
$\tau=K|L-J|/2\le K|I|/2=K/18$。
因此远区 $1\le\tau\le K/18$ 的原始网格距离 $q$ 数为
$O_{T,A,D}(Kv)$，其归一误差至多
\[
 O_{p,w,T,A,D}(Kv)e^{-v/25}.
\]
由 $v\ge16\log n$ 和 $Kv=c_p\sqrt{nv}=O_p(\sqrt{n\log n})$，
\[
 Kv e^{-v/25}
 \le C_p\sqrt{n\log n}\,n^{-16/25}
 =C_p\sqrt{\log n}\,n^{-7/50}\longrightarrow0.
\]
将区间端点任意分配到相邻一区，可使近、中、远三区互斥；三者又穷尽
所有不同节点长度。对节点交集及段对求和，合并三个估计即得结论。
全过程没有假设段事件独立，命题中的比较误差始终乘实际亲和度
$H_{LJ}^n$，中、远距离的衰减因而同样保留。
\end{proof}

\subsection{有限质量截断与主定理}
\begin{proof}[定理\ref{thm:main}的证明]
先固定 $p,T,A,N,D,w$ 如上，并记
\[
 G_w(L)=E\exp\!\left(\max_{0\le a\le N}
 X_{1/w,\lambda(L)}(aA)\right),\qquad
 \lambda(L)=\frac{58}{9L^2}.
\]
有限节点最大值等式\eqref{eq:paperwinner}给
\[
 ue^hP(U_i)=\phi(0)G_w(L_i)+\varepsilon_{i,n},
 \qquad\sup_{i\in\mathcal I_n}|\varepsilon_{i,n}|\longrightarrow0.
\]
这里每个 $U_i$ 的实际长度、实际 $q_{L}$ 和自己的倾斜都已计入该等式，
并非另假设各段同分布或方差等于 $v$。

函数 $G_w$ 在闭区间 $I$ 上连续有界。事实上有限 高斯 向量的
均值和协方差随 $\lambda(L)$ 连续，故有限向量弱收敛；而
\[
 E e^{2\max_aX_{1/w,\lambda}(aA)}
 \le\sum_{a=0}^NEe^{2X_{1/w,\lambda}(aA)}
 \le(N+1)e^{\lambda_{\max}T}
\]
给出紧参数区间上的一致可积性，故指数最大值的期望也连续。
原基点网格间距为 $\Delta_n=2P_0/(Kv)\to0$，末节点条件只排除
$I$ 右端长度 $2T/(Kv)$ 的小条带。因此原完整基点和满足
\[
 \Delta_n\sum_{i\text{原基点}}G_w(L_i)\longrightarrow\int_I G_w(L)\,dL.
\]
安全删除了 $o(Kv)$ 个基点；因 $G_w$ 有界，所删和乘 $\Delta_n$
的绝对值为 $o(1)$。所以同一 Riemann 极限对 $\mathcal I_n$ 成立。
又统一余项的总贡献满足
\[
 \frac1{u\sqrt n}\left|\sum_{i\in\mathcal I_n}\varepsilon_{i,n}\right|
 \le\frac{J_*}{u\sqrt n}\sup_i|\varepsilon_{i,n}|\longrightarrow0.
\]
由于 $Kv/(u\sqrt n)=c_p$，遂得完整一阶极限
\begin{equation}\label{eq:paperfirstsum}
 \frac{e^h}{\sqrt n}\sum_{i\in\mathcal I_n}P(U_i)
 \longrightarrow\frac{c_p\phi(0)}{2P_0}\int_I G_w(L)\,dL.
\end{equation}
上一引理与段数渐近又给
\begin{equation}\label{eq:papersecondsum}
 \limsup_n\frac{e^h}{\sqrt n}
 \sum_{\substack{i,j\in\mathcal I_n\\i<j}}P(U_i\cap U_j)
 \le\frac{c_p|I|\phi(0)}{2P_0}e^{4/(3w)}(N+1)^2
 \frac{e^{-cD/4}}{1-e^{-cP_0/4}}.
\end{equation}
在每个有限维数，先施二阶 Bonferroni 不等式：
\[
 P\left(\bigcup_{i\in\mathcal I_n}U_i\right)
 \ge\sum_iP(U_i)-\sum_{i<j}P(U_i\cap U_j).
\]
所有 $S_{L_{i,a}}$ 都是有限个可测函数的和，每段只有 $N+1$ 个节点，
故 $U_i$ 可测。每个固定 $n$ 的段数 $J_*$ 又是有限的；上述并集与二阶
Bonferroni 均在有限集合上应用，没有无限并事件的交换问题。
利用式\eqref{eq:paperfirstsum}的一阶极限及式\eqref{eq:papersecondsum}
的二阶上界，可以在随后取 $n$ 极限时保留二者之差。

现在将周期构造变成 $L^1(\mathbb R^n)$ 的实际见证。记
$F_n(x)=\prod_{a=1}^n\rho(x_a)$ 为周期乘积。对每个共同边长 $L$，
Fubini 定理使中心立方体平均准确分解为
\[
 L^{-n}\int_{[-L/2,L/2]^n}F_n(x+y)\,dy
 =\prod_{a=1}^n\left(\frac1L\int_{-L/2}^{L/2}\rho(x_a+y)\,dy\right)
 =\prod_{a=1}^nY_L(x_a)=e^{S_L(x)}.
\]
因此在 $\bigcup_iU_i$ 上，至少一个实际中心立方体平均严格超过 $e^h$。
不同尺度使用同一组坐标；乘积公式没有把尺度之间的事件变成独立事件。

令整数 $R=n^2$，定义
\[
 f_n(x)=\mathbf1_{[-R,R]^n}(x)F_n(x).
\]
一维因子严格正、有界且周期质量1，所以 $f_n$ 非负、有界、紧支撑，且
\[
 \|f_n\|_1=\prod_{a=1}^n\int_{-R}^R\rho(x_a)\,dx_a=(2R)^n.
\]
对内盒 $[-R+1,R-1]^n$ 内的中心，由于全部所用 $L\in I\subset(1,2)$，
各中心立方体的半边长小于1，因而完全包含于外盒 $[-R,R]^n$；
这些平均与 $F_n$ 的平均完全相同。
事件 $\bigcup_iU_i$ 对每个坐标都是周期1的可测集合。内盒边长
$2R-2$ 是整数，按单位周期盒分解，其事件体积恰为
\[
 (2R-2)^nP^{\otimes n}\left(\bigcup_iU_i\right).
\]
因此在弱型比率中取阈值 $e^h$，得到
\begin{align*}
 C_n&\ge\frac{e^h|\{M_nf_n>e^h\}|}{\|f_n\|_1}\\
 &\ge e^h\left(1-\frac1R\right)^n
 P^{\otimes n}\left(\bigcup_iU_i\right).
\end{align*}
此处极大函数的可测性可独立于任何弱型结论核验。对任意
$f\in L^1(\mathbb R^n)$ 和固定中心 $x$，当 $L\to L_0>0$ 时，
两个中心立方体的对称差体积趋于0，Lebesgue 积分的绝对连续性给
其上 $|f|$ 的积分趋于0；再用 $L^{-n}\to L_0^{-n}$，得到
$L\mapsto A_L|f|(x)$ 连续。因此 $M_nf(x)$ 的尺度上确界可以限制为
正有理 $L$ 的可数上确界。每个 $A_L|f|$ 由可测卷积积分给出，
所以 $M_nf$ 及上述严格超水平集均可测。该论证没有使用待证的弱型估计。
Bernoulli 不等式给
\[
 \left(1-\frac1{n^2}\right)^n\ge1-\frac1n\longrightarrow1.
\]
故此截断损失不改变归一化的极限下界。结合 Bonferroni 和上述两个和式，
对每组固定的有限参数有
\begin{align*}
 \liminf_{n\to\infty}\frac{C_n}{\sqrt n}
 \ge\frac{c_p\phi(0)}{2(T+D)}\left[
 \int_I G_w(L)\,dL
 -|I|e^{4/(3w)}(N+1)^2
 \frac{e^{-cD/4}}{1-e^{-c(T+D)/4}}\right].
\end{align*}

保持 $p,T,A,N,D$ 固定，下一步仅在已经成立的数值下界中令
$w\to\infty$。有限节点下，箱核
$b_{[0,t]}=\mathbf1_{[0,t]}*\chi_{1/w}$ 在 $L^2$ 中趋于
$\mathbf1_{[0,t]}$；前述有限向量收敛及一致可积性于是给
\[
 G_w(L)\longrightarrow
 E\exp\!\left(\max_{0\le a\le N}
 \{\sqrt{\lambda(L)}B(aA)-\lambda(L)aA/2\}\right)
 =H_{\kappa(L)}(T,A)
\]
在 $L\in I$ 上一致，其中
$\kappa(L)=\lambda(L)/2=29/(9L^2)$。
一致性也可由紧参数子列论证得到：任意失败的基点子列有
$\lambda(L)$ 收敛子列，而同一有限 高斯 收敛及二阶指数矩界排除失败。
同时 $e^{4/(3w)}\to1$。因而
\begin{equation}\label{eq:paperfinitebound}
 \liminf_n\frac{C_n}{\sqrt n}\ge
 \frac{c_p\phi(0)}{2(T+D)}
 \left[\int_I H_{\kappa(L)}(T,A)\,dL
 -|I|(N+1)^2\frac{e^{-cD/4}}{1-e^{-c(T+D)/4}}\right].
\end{equation}
这里 $w$ 趋于无穷没有改变先前的维数极限：该不等式来自对所有有限
$w$ 已成立的右端取极限，左端始终是同一个数值下界。

现对每个整数 $T\ge2$ 选择
\[
 A=T^{-2},\qquad N=T^3,\qquad D=(32/c)\log T.
\]
对每个这样的 $T$，$A,N,D$ 仍是有限固定参数；为了使用有限节点结论，
先取满足 $w\ge\max(4,T^2)$ 的有限 $w$，先取全部整数 $n$ 极限，
再取 $w$ 的上确界，才得到式\eqref{eq:paperfinitebound}。
随后对这些已经建立的数值下界令整数 $T\to\infty$。有
$D/T=(32/c)(\log T)/T\to0$，而
\begin{align*}
 (N+1)^2\frac{e^{-cD/4}}{1-e^{-c(T+D)/4}}
 &=(T^3+1)^2\frac{T^{-8}}{1-e^{-c(T+D)/4}}
 =O(T^{-2}).
\end{align*}
分母趋于1；特别是它最终有一个正的统一下界。
网格引理\ref{lem:papermesh}又给
\[
 \sup_{L\in I}\left|\frac{H_{\kappa(L)}(T,T^{-2})}T-\kappa(L)\right|
 \longrightarrow0,
\]
故可以在固定有限区间 $I$ 上积分。
把式\eqref{eq:paperfinitebound}的分子、分母同除以 $T$，交集项除以
$T$ 后为 $O(T^{-3})$，于是对每个固定 $p<1$ 得
\begin{align*}
 \liminf_n\frac{C_n}{\sqrt n}
 &\ge\frac{c_p\phi(0)}2\int_{10/9}^{11/9}\frac{29}{9L^2}\,dL\\
 &=\frac{(3p/\sqrt{29})(1/\sqrt{2\pi})}{2}
 \frac{29}{9}\left(\frac9{10}-\frac9{11}\right)\\
 &=\frac{(3p/\sqrt{29})(1/\sqrt{2\pi})}{2}\frac{29}{110}
 =p\frac{87}{220\sqrt{58\pi}}.
\end{align*}
最后对所有固定 $0<p<1$ 已成立的数值下界取上确界，得到
\[
 \liminf_n\frac{C_n}{\sqrt n}\ge\frac{87}{220\sqrt{58\pi}}.
\]
这一步没有把 $p=1$ 代入原密度，严格正的背景始终保留在每个实际有限参数见证中。
\end{proof}

\begin{corollary}[统一维数与见证的正则性]\label{cor:paperwitness}
对任意 $b<87/(220\sqrt{58\pi})$，存在 $n_b$，使每个整数 $n\ge n_b$
都有 $C_n\ge b\sqrt n$。每个这样的维数可用非负、有界、紧支撑的
乘积截断函数见证此下界；其一维周期因子严格正、质量为1，且高度为
$O_b(\log n)$。
\end{corollary}
\begin{proof}
该见证断言需要保留严格余量，不能仅由最佳常数的上确界形式断言它被某个函数取得。
记 $C_*=87/(220\sqrt{58\pi})$，选 $b_1$ 使
$\max(b,0)<b_1<C_*$，再选固定 $p<1$ 使 $pC_*>b_1$。
上一证明给出的有限 $T$ 右端随 $T\to\infty$ 趋于 $pC_*$，
故可选一个有限整数 $T$，以及相应有限的 $A=T^{-2}$、$N=T^3$、
$D=(32/c)\log T$，使式\eqref{eq:paperfinitebound}的右端严格大于 $b_1$。
对这组固定参数，有限箱宽右端随 $w\to\infty$ 收敛到该右端；
因此又可选一个有限 $w\ge\max(4,T^2)$，使尚未令 $w\to\infty$ 的
固定参数下界仍严格大于 $b_1$。

对这些最终固定的 $p,T,A,N,D,w$，实际截断见证的弱型比率记为
\[
 \mathcal R_n=\frac{e^h|\{M_nf_n>e^h\}|}{\|f_n\|_1\sqrt n}.
\]
上述一阶极限、二阶上界和体积损失共同给
$\liminf_n\mathcal R_n>b_1>b$。
由下极限的定义，存在参数依赖的统一阈值，之后每个整数 $n$ 都有
$\mathcal R_n>b$；因 $C_n/\sqrt n\ge\mathcal R_n$，这就给出所需实际见证。

说明这一阈值确实覆盖全部充分大整数维数。对每个整数 $n$，增长处方
选使 $v=9p^2n/(29K^2)\ge16\log n$ 成立的最大 $K=2^d-1$。
下一允许周期为 $2K+1$，其不再满足此不等式，故
\[
 16\log n\le v<(2+1/K)^2\,16\log n.
\]
由此 $v=\Theta_p(\log n)$、$K\asymp_p\sqrt{n/\log n}$，
且 $K/v^r\to\infty$ 对每个固定 $r$ 成立。最小二幂 $B\ge v$
满足 $B<2v$，所以有限域存在阈值 $K+1>H(B)$、
完整核的充分条件、正性条件及所有固定参数概率比较条件，均在某个
统一阈值后成立。周期切换两侧仍满足同一两边估计，不曾抽取某个维数子序列。
取这些有限阈值与 $\mathcal R_n>b$ 的阈值之最大值，即得 $n_b$。

最后，原密度的解析上下界给
\[
 \rho(x)\ge1-p-\frac{20pH_M}{3K},\qquad
 \rho(x)\le1-p+pwv+\frac{20pH_M}{3K},
 \quad H_M=\sum_{m\le M}\frac1m.
\]
因为 $p<1$ 固定，$H_M/K\to0$，所以最终下界至少为 $(1-p)/2>0$。
又 $v=O_p(\log n)$、$w$ 固定，故一维高度为 $O_b(\log n)$；
周期质量恰为1来自每峰质量 $1/K$ 和有限低频项的零均值。
乘积在每个固定 $n$ 中有界，乘以有限盒指示函数后紧支撑，因而属于
$L^1(\mathbb R^n)$。这里 $O_b(\log n)$ 指一维因子高度，整个
$n$ 维乘积的高度不被宣称具有同一量级。
\end{proof}

\begin{remark}[结论的量词]
合法顺序是先固定有限参数并取全部整数 $n\to\infty$，再在数值下界中
令 $w\to\infty$，再令整数 $T\to\infty$，最后取 $p\uparrow1$ 的上确界。
每个严格较小系数由一组有限参数及其实际截断见证实现。
此证明没有提供 $T,N,w$ 同时随 $n$ 增长的统一误差率，也不提供显式 $n_b$。
\end{remark}

\section{常数与结论的范围}\label{sec:discussion}
\begin{proof}[定理\ref{thm:rational}的证明]
由 $\pi<22/7$ 及
\[
 (68/5)^2-58\cdot\frac{22}{7}=\frac{468}{175}>0
\]
得
\[
 C_*>\frac{87}{2992}=\frac1{35}+\frac{53}{104720}>\frac1{35}.
\]
定理\ref{thm:main}和极限下确界的定义即给出结论。
\end{proof}

常数 $C_*$ 来自局部 Brownian 越阈强度在尺度区间 $I=[10/9,11/9]$ 上的积分，
而非有限样本回归。具体地，强度 $\kappa(L)=29/(9L^2)$ 给出
\[
 \int_{10/9}^{11/9}\kappa(L)\,dL=\frac{29}{110},\qquad
 \frac{3}{2\sqrt{29}}\phi(0)\frac{29}{110}=C_*.
\]
数值为 $C_*\approx0.0292959535703970$。这一小数仅用于比较常数大小；
上面的有理不等式已严格证明其大于 $1/35$。

本构造的每个测试函数均属于 $L^1\cap L^\infty$，但其上确界范数随维数增长。
参数依次取极限，使背景项在构造时保持严格正性，也使每次高斯比较只涉及固定的
有限维向量。证明不提供一个可直接计算的充分维数，也不推出所有低维均满足
$C_n\ge\sqrt n/35$。附带数值报告据此将确定性算术检查、有限参数实验和
渐近命题分别列出。

结合定理 A，本文得到的增长区间为
\[
 \sqrt n\lesssim C_n\lesssim n\log(e+n).
\]
是否存在同为平方根量级的上界，或下界还能超过平方根阶，均不由本文解决。

\begin{thebibliography}{99}
\bibitem{SS1983} E. M. Stein and J.-O. Strömberg.
Behavior of maximal functions in $\mathbb R^n$ for large $n$.
\emph{Arkiv för Matematik}, 21 (1983), 259--269.
\href{https://doi.org/10.1007/BF02384314}{doi:10.1007/BF02384314}.
\bibitem{Aldaz2011} J. M. Aldaz.
The weak type $(1,1)$ bounds for the maximal function associated to cubes grow to infinity with the dimension.
\emph{Annals of Mathematics}, 173 (2011), 1013--1023.
\href{https://doi.org/10.4007/annals.2011.173.2.10}{doi:10.4007/annals.2011.173.2.10}.
\bibitem{Aubrun2009} G. Aubrun.
Maximal inequality for high-dimensional cubes.
\emph{Confluentes Mathematici}, 1 (2009), no.~2, 169--179.
\href{https://doi.org/10.1142/S1793744209000067}{doi:10.1142/S1793744209000067}.
\bibitem{IS2013} A. S. Iakovlev and J.-O. Strömberg.
Lower bounds for the weak type $(1,1)$ estimate for the maximal function associated to cubes in high dimensions.
\emph{Mathematical Research Letters}, 20 (2013), no.~5, 907--918.
\href{https://doi.org/10.4310/MRL.2013.v20.n5.a7}{doi:10.4310/MRL.2013.v20.n5.a7}.
\bibitem{Bourgain2012} J. Bourgain.
On the Hardy--Littlewood maximal function for the cube.
Preprint, 2012. \href{https://arxiv.org/abs/1212.2661}{arXiv:1212.2661}.
\bibitem{Raic2019} M. Raič.
A multivariate Berry--Esseen theorem with explicit constants.
\emph{Bernoulli}, 25 (2019), no.~4A, 2824--2853.
\href{https://doi.org/10.3150/18-BEJ1072}{doi:10.3150/18-BEJ1072}.
\bibitem{Shevtsova2011} I. Shevtsova.
On the absolute constants in the Berry--Esseen type inequalities for identically distributed summands.
Preprint, 2011. \href{https://arxiv.org/abs/1111.6554}{arXiv:1111.6554}.
\end{thebibliography}

\end{document}
